\documentclass[11pt]{amsart}
\usepackage[T1]{fontenc}
\usepackage{lmodern,amsmath,amssymb,amsthm,mathtools,microtype,booktabs,array,longtable}
\usepackage[margin=1.06in]{geometry}
\usepackage{tikz}
\usepackage{xr-hyper}
\usepackage[hidelinks]{hyperref}
\hypersetup{pdftitle={Finite-Use Geometry and Learning of Ordered Quantum Measurements},pdfauthor={Qian Qi}}
\emergencystretch=2em
\numberwithin{equation}{section}
\newtheorem{theorem}{Theorem}[section]
\newtheorem{lemma}[theorem]{Lemma}
\newtheorem{proposition}[theorem]{Proposition}
\newtheorem{corollary}[theorem]{Corollary}
\theoremstyle{definition}\newtheorem{definition}[theorem]{Definition}
\theoremstyle{remark}\newtheorem{remark}[theorem]{Remark}
\newcommand{\Q}{\mathbb Q}\newcommand{\R}{\mathbb R}\newcommand{\C}{\mathbb C}\newcommand{\N}{\mathbb N}\newcommand{\Z}{\mathbb Z}\newcommand{\E}{\mathbb E}
\newcommand{\ip}[2]{\langle #1,#2\rangle}\newcommand{\norm}[1]{\lVert #1\rVert}
\newcommand{\epsc}{\varepsilon_{\mathrm c}}
\DeclareMathOperator{\osc}{osc}\DeclareMathOperator{\spanop}{span}\DeclareMathOperator{\End}{End}\DeclareMathOperator{\diag}{diag}\DeclareMathOperator{\tr}{tr}\DeclareMathOperator{\rad}{rad}\DeclareMathOperator{\TV}{TV}
\title[Finite-use measurement geometry and learning]{Finite-Use Geometry and Learning of Ordered Quantum Measurements}
\author{Qian Qi}
\date{5 October 2026}
\subjclass[2020]{Primary 81P15, 81P47; Secondary 94A17, 81P70}
\externaldocument{build/xref-supplement}[BINARY_SUPPLEMENT.pdf]
\begin{document}
\raggedbottom
\begin{abstract}
We study finite-use discrimination of ordered quantum measurements
with classical output.  At a fixed measurement $E$ and nonzero
one-sided tangent $H$, we classify all legal finite perturbations
$F=E+sH+O(s^2)$ with a prescribed operator-norm remainder bound.
The adaptive trace separation has order $\min\{1,\sqrt N s\}$ when
$H$ belongs to the normalized covariance range and
$\min\{1,Ns\}$ otherwise, uniformly over that remainder neighborhood.
The range criterion is expressed by the missing-support blocks and
the support generator $\Gamma_E(H)=\frac i2\sum_j[\operatorname{supp}E_j,H_j]$.
Linear separation has two mechanisms: a first-order support opening
is already detected by independent probes, whereas a tangential
coherent direction has only square-root separation under product
probe--reference acquisition, even with collective final readout.
The results include every fixed one-sided $C^2$ curve with nonzero
first derivative, and higher-order approaches with a quadratic
remainder in the leading parameter.  An exact scalar family has the
mixed scale $\min\{1,\sqrt N t^a+Nt^b\}$ for integers $1\le a<b<2a$.
The proof combines the complete support kernel with finite normalized
surrogates, impossible events, and the established metrological
correction mechanism.  Constants depend on the fixed base, direction
and remainder bound; the global arbitrary-pair covariance result
remains an upper certificate.  Exact cone and supplied-pair decisions
are polynomial on rational input.  The existing transported
regularization and balanced-interior coding and learning theorems
are retained, with fixed-outcome sharpness and trusted-control
resources stated separately.
\end{abstract}
\maketitle
\section{Introduction}\label{sec:introduction84}
A measurement with a classical output may retain a quantum advantage
when its input is correlated with a reference.  The question considered
here is quantitative: how well can two ordered measurements be
separated with at most $N$ uses, and which features of their effects
determine the dependence on $N$?  The answer must accommodate both
strictly positive effects and singular measurements.  Positivity of
each effect alone is not the relevant constraint: the effects sum to
the identity, and their boundary geometry is coupled by this equation.

\subsection{The model and the principal results}
An ordered $k$-outcome measurement on $\C^d$ is a tuple
$E_j\succeq0$ with $\sum_jE_j=I_d$.  One call is the channel
\[
 \mathcal M_E(\rho)=\sum_{j=1}^k\tr(E_j\rho)|j\rangle\langle j|.
\]
The input is consumed and the output is only the classical label;
no residual device system is available.  A tester may prepare inputs
entangled with its own reference, retain quantum memory, and choose
subsequent controls using previous labels.  It uses the same control
rule under both hypotheses.  Its public stopping time is bounded
by $N$.  We define $D_N(E,F)$ to be the supremum of the unhalved
trace distance between its two final states.  Thus $0\le D_N\le2$.
Nonadaptive testers may have entangled inputs and retained references;
``nonadaptive'' does not mean ``product input.''

\subsection{A finite neighborhood theorem and three acquisition regimes}
Fix a base measurement $E$ and a nonzero Hermitian direction $H$ with
$\sum_jH_j=0$.  Write $P_j=\operatorname{supp}E_j$, $Q_j=I-P_j$, and
\[
 J_j=Q_jH_jQ_j,\qquad
 \Gamma_E(H)=\frac i2\sum_j[P_j,H_j],\qquad
 \mathcal W_E=\sum_jP_j\mathcal H_dP_j.
\]
The exact cone of right derivatives of analytic measurement curves is
$J_j\succeq0$ for all $j$.  It includes first-order support openings;
its lineality space is $J_j=0$.  A rational normalized matrix curve
proves sufficiency even when a nonzero $J_j$ is singular and has
coupled cross-support blocks (Lemma~\ref{lem:onesidedcone86}).

Our finite-pair theorem fixes this $E,H$ and a remainder bound
$\Lambda\ge0$.  It allows \emph{every} legal $F$ satisfying
\[
 \sum_j\|F_j-E_j-sH_j\|_{\mathrm{op}}\le\Lambda s^2.
\]
No curve is specified, and the remainder may vary arbitrarily with
$s$.  Theorem~\ref{thm:tubedichotomy86} proves, uniformly in all such
$F$, every $N\ge1$ and all sufficiently small $s>0$,
\[
 D_N(E,F)\asymp_{E,H,\Lambda}
 \begin{cases}
 \min\{1,\sqrt N s\},&J_j=0\ \hbox{for all }j,
                         \ \Gamma_E(H)\in\mathcal W_E,\\
 \min\{1,Ns\},&\hbox{otherwise}.
 \end{cases}
\]
These constants are uniform over the finite remainder neighborhood,
not over the base point or the direction.  This does not assert
arbitrary-pair equivalence of the global covariance certificate.

There are two distinct linear mechanisms.  A nonzero $J_j$ creates an
outcome impossible at the base; repeated product inputs then suffice.
If every $J_j=0$ but $\Gamma_E(H)\notin\mathcal W_E$, a corrected
logical qubit gives the linear lower.  For product probe--reference
inputs, with no inter-pair entanglement or feedback but arbitrary
collective processing of completed outputs, the latter regime has
only the square-root order.  Theorem~\ref{thm:producttrichotomy86}
proves all three regimes, including a reference-uniform converse
for that acquisition class.  Arbitrary entangled parallel acquisition
is a different class and is not bounded by that converse.

The Kraus-span metrological exponents and correction mechanism are
established principles \cite{ZJ84,KL84}.  The results here give the
complete effect-coordinate cone, a finite-neighborhood comparison,
and the distinction between coherent accumulation and support-opening
accumulation.  Their proofs retain all finite remainder terms and do
not infer trace-distance bounds from an asymptotic Fisher-information
statement.  Product fidelity and its tensorization are standard
\cite[Chapter~3]{Watrous65}.

The neighborhood theorem includes one-sided $C^2$ curves with nonzero
right derivative (Corollary~\ref{cor:onesidedcurves86}), and higher-order
families $E+t^qH+O(t^{2q})$ (Corollary~\ref{cor:powerjets86}).
A scalar three-outcome family with an intervening support opening has
instead the mixed law $\min\{1,\sqrt N t^a+Nt^b\}$ for
$1\le a<b<2a$ (Proposition~\ref{prop:mixedrates86}).  Thus a zero
first derivative alone remains insufficient for classification.
The result does not silently assume away higher-order leakage.

\subsection{The covariance and support structure}
Our complete-body upper bound uses a positive, concave covariance
$\mathcal C_E$ on zero-sum Hermitian tuples.  Its factorization through
a horizontal projection gives a finite-use bound valid through changes
of rank and zero effects.  With $M=(E+F)/2$ and $H=F-E$,
\begin{equation}\label{eq:frontupper84}
 D_N(E,F)\le\min\left\{2,2\sqrt{k+1}
  \sqrt{N\langle H,(\mathcal C_M+N^{-1}\operatorname{Id})^{-1}H\rangle}
                              \right\}.
\end{equation}
This is a global upper certificate, not a claimed formula for the
optimal adaptive distance.  The normalization includes the exact
binary Sylvester restriction of Theorem~\ref{thm:closedcovariance83}.

The support description determines its entire kernel.  If
$P_j=\operatorname{supp}E_j$ and
$\mathcal W_E=\sum_jP_j\mathcal H_dP_j$, then
Theorem~\ref{thm:supportkernel84} parametrizes all null directions
and proves
\[
 \dim\ker(\mathcal C_E|_{\mathcal T})
 =d^2-\dim\mathcal W_E+\sum_j(d-\operatorname{rank}E_j)^2.
\]
This distinction between support span and the mere effect span is
essential for effects of rank greater than one.  In particular,
a nonprojective measurement can have a nontrivial covariance kernel;
the earlier projective characterization concerns a zero operator,
not one null direction.

This kernel has a finite-use interpretation on every fixed unitary
orbit $E_j(t)=e^{itG}E_je^{-itG}$.  Apart from the stationary case,
Theorem~\ref{thm:orbittrichotomy84} proves, uniformly in every integer
$N\ge1$ and all sufficiently small $|t|$,
\begin{equation}\label{eq:frontorbit84}
 D_N(E,E(t))\asymp_{E,G}
 \begin{cases}
 \min\{1,\sqrt N|t|\},&G\in\mathcal W_E,\\
 \min\{1,N|t|\},&G\notin\mathcal W_E.
 \end{cases}
\end{equation}
The second lower bound uses a corrected logical qubit and the explicit
finite-angle estimate \eqref{eq:finiteanglelower84}.
The support span is precisely the Hermitian Kraus-product span of
this measurement channel.  The metrological criterion is credited
to Zhou--Jiang \cite{ZJ84}; the present formulation identifies it
with the covariance kernel and proves the finite-pair trace estimates
directly.  This is an orbit theorem on the complete body, not a
full-boundary volume or minimax theorem.

We also replace a coordinate-fixed ridge by a public scalar-measurement
covariance.  Transporting its probability vector through the same
output channel gives a regularized data-processing inequality and
exact invariance under reversible label splitting
(Theorem~\ref{thm:transportedridge84}).  The original fixed Euclidean
ridge is not asserted to have this property.

\subsection{Nonunitary curves and certified controls}
The support description also gives a local classification without a
unitary-orbit assumption.  For a fixed two-sided $C^2$ measurement curve
with $H=E'(0)\ne0$, set
\[
 \Gamma_E(H)=\frac i2\sum_j[\operatorname{supp}E_j,H_j].
\]
Theorem~\ref{thm:curveclassification85} proves
\[
 D_N(E(0),E(t))\asymp_{E(\cdot)}
 \begin{cases}
 \min\{1,\sqrt N|t|\},&\Gamma_E(H)\in\mathcal W_E,\\
 \min\{1,N|t|\},&\Gamma_E(H)\notin\mathcal W_E.
 \end{cases}
\]
This holds for every finite $N$ on a two-sided interval depending on
the fixed curve.  The supports may change at second order.  An analytic
horizontal surrogate supplies a $\sqrt N|t|+O(Nt^2)$ upper in the
range branch; a support-complement code supplies a logical channel
with an explicit $O(t^2)$ remainder in the other branch.  Both are
finite inequalities rather than inferences from asymptotic Fisher
information.  The metrological exponents and correction mechanism
remain attributed to \cite{ZJ84,KL84}; the new formulation works
directly with $C^2$ effects even when an exact Kraus family is not
assumed differentiable at a rank opening.

Theorem~\ref{thm:controlstability85} bounds the loss from uniformly
certified errors in preparation, encoding, recovery and readout.
The allowed per-cycle error scales with $|t|\Delta$, where $\Delta$
is the logical gap.  This is a conditional stability guarantee for
known-pair controls, not an executed implementation or a gate-synthesis
complexity theorem.  The constants and interval may degenerate with
support eigenvalues, tangent size, curvature and the support residual.
Zero first derivative does not imply stationarity for a general curve.
The global arbitrary-pair covariance bound remains one-sided.

\subsection{Learning and descriptions on balanced interiors}
The operational consequences are complemented by exact description
and common learning laws on the balanced interior
$I_d/(2k)\preceq E_j\preceq3I_d/(2k)$.  Its real affine dimension
is $p=(k-1)d^2$.  Theorems~\ref{thm:multioutcomeentropy82} and
\ref{thm:multioutcomecode82} give
\[
 \log_2\mathcal C_{N,k}(\delta)
 =\frac p2\log_2N+p\log_2(1/\delta)+O(p\log(k+1)),
 \qquad 0<\delta<1/128.
\]
The common learner has the joint order
\[
 M^*_{d,k}(N,\delta,\eta)
 =\Theta_k\!\left(N\delta^{-2}[d^2+d\log(1/\eta)]\right)
\]
for fixed $k$, $d,N\ge1$, $0<\delta\le2^{-24}$ and
$0<\eta\le1/8$.  The displayed constructive upper has a $k^3$
factor; the lower does not match growing $k$.
Theorem~\ref{thm:affinerepair83} supplies polynomial, search-free
statistical legalization, while the entropy-optimal public dictionary
remains a separate finite, potentially exhaustive construction.
The lower theorem permits arbitrary coherent adaptive training,
whereas the upper uses fresh acquisitions and collective processing
of completed outputs.  Pair-dependent discrimination witnesses do
not serve as a common learning design.

\subsection{Resource conventions and reading order}
Unknown-device calls, future-use horizon, transmitted index length,
public dictionary reconstruction, classical arithmetic, trusted-control
synthesis, and physical execution are distinct resources.  A query
bound here holds on every record, with a defined legal fallback.
Public dictionaries and their parameters are independent of the
target; target-dependent private advice would count as transmitted
information.  Exact rational certificate evaluation and balanced
legalization have polynomial bit complexity.  Neither statement
asserts polynomial synthesis of an arbitrary collective readout or
of the whole entropy-optimal dictionary.  The orbit recovery is an
ideal-control comparison construction, not an executed common learner.

Sections~\ref{sec:covariance83}--\ref{sec:ridge84} develop the global
upper, support kernel, orbit classification and transported ridge.
Sections~\ref{sec:finiteoutcome82}--\ref{sec:affinerepair83} give the
balanced-interior learning and coding consequences.
All binary auxiliary proofs used by this article are supplied in the
separate \emph{Binary Foundations and Auxiliary Proofs} supplement.
Cross-document theorem references carry the prefix S and are resolved
from its current source, not from a historical PDF.  The primary article
and supplement together are a self-contained proof package; the
independent structural paper is not a premise of any measurement theorem.

\subsection{Logical prerequisites of the primary article}
The covariance factorization, variance identity and full-body path
upper are proved in Section~\ref{sec:covariance83}; the finite
Bernoulli lower is proved in Section~\ref{sec:product85}.  Together
with the interface above, these suffice for the support and
finite-discrimination results of
Sections~\ref{sec:support84}--\ref{sec:controlstability85}.
The independently complete structural paper is never a premise.
The finite-neighborhood and one-sided results of
Sections~\ref{sec:tangentneighborhood86}--\ref{sec:independent86}
use these same premises; the product converse adds only standard
fidelity properties stated in its proof.
\begin{center}
\begin{tabular}{@{}>{\raggedright\arraybackslash}p{.31\textwidth}>{\raggedright\arraybackslash}p{.61\textwidth}@{}}
\toprule
Primary result & Required auxiliary result\\
\midrule
Support kernel, orbit, curve and finite-neighborhood classification &
None from the supplement: all probability, covariance and correction
premises are proved in the primary.\\
Independent product probes &
Primary normalized-factor remainder and the explicitly stated
fidelity purification, tensorization and trace inequalities.\\
Transported regularization &
Primary covariance Jensen identity and variational duality.\\
Balanced-interior common learning &
Supplement Theorem~\ref{thm:rationallearner81} and
Corollary~\ref{cor:operatorconfidence80}; the former imports the
credited binary tomography primitive.\\
Noisy-projector comparison &
Supplement Theorem~\ref{thm:matrixmetric76} for the specified binary
pair. It is not a premise of the smooth-curve theorem.\\
Other binary boundary, coding and control proofs &
Preserved in the current supplement; not premises of the new
support or smooth-curve classification.\\
\bottomrule
\end{tabular}
\end{center}

\section{A covariance principle on the closed measurement body}
\label{sec:covariance83}

Let $\mathcal P_{d,k}$ denote the complete body of ordered $k$-outcome
measurements on $\C^d$: its elements are Hermitian tuples
$E=(E_1,\ldots,E_k)$ with $E_j\succeq0$ and $\sum_jE_j=I_d$.
Here $d\ge1$, $k\ge2$, and $N\ge1$ are integers.  Zero effects,
changing ranks, and noncommuting supports are allowed throughout
this section.  The operational distance $D_N$ uses unhalved trace
norm, retained references, common adaptive control, and at most $N$
calls with bounded public stopping.  The device consumes its input
and returns only its ordered classical label.

The normalized tangent space and its inner product are
\[
 \mathcal T_{d,k}=\{H_j=H_j^*: \textstyle\sum_jH_j=0\},\qquad
 \langle H,K\rangle=\sum_j\tr(H_jK_j).
\]
All operator inequalities below on tuple spaces refer to this real
Hilbert structure, unless stated otherwise.  For any Hermitian tuple
$K$, not necessarily of zero sum, put
\begin{align}
 S_E(K)&=\sum_jE_jK_j,\label{eq:covariances83}\\
 (\mathcal C_EK)_j
  &=\tfrac12\bigl(E_jK_j+K_jE_j-E_jS_E(K)-S_E(K)^*E_j\bigr).
                                      \label{eq:covarianceoperator83}
\end{align}
The matrix $S_E(K)$ need not be Hermitian.  Keeping its adjoint in
\eqref{eq:covarianceoperator83} is essential.  The operator
$\mathcal C_E$ is polynomial in the effects, annihilates constant
tuples, and has range in $\mathcal T_{d,k}$.  In particular its
restriction to $\mathcal T_{d,k}$ is well defined.

For $\tau=N^{-1}$ set
\begin{align}
 g_{N,E}(H)^2
  &=N\langle H,(\mathcal C_E+\tau\operatorname{Id})^{-1}H\rangle,
                                      \label{eq:covarianceform83}\\
 \mathcal Q_N(E,F)&=g_{N,(E+F)/2}(F-E).
                                      \label{eq:covariancemidpoint83}
\end{align}
The inverse acts on $\mathcal T_{d,k}$.  Positivity proved below
makes it nonsingular even on the spectral boundary.  We do not
identify $\mathcal Q_N$ with an operational distance or assume a
triangle inequality for it.

\begin{theorem}[Closed-body adaptive bound]
\label{thm:closedcovariance83}
For every $E,F\in\mathcal P_{d,k}$,
\begin{equation}\label{eq:closedcovariance83}
 D_N(E,F)\le
 \min\{2,\,2\sqrt{k+1}\,\mathcal Q_N(E,F)\}.
\end{equation}
More generally, every continuously differentiable path $E(t)$ in
$\mathcal P_{d,k}$ satisfies
\begin{equation}\label{eq:covariancepath83}
 D_N(E(0),E(1))\le
 \sqrt{k+1}\int_0^1g_{N,E(t)}(E'(t))\,dt.
\end{equation}
These statements are uniform over all ranks and support patterns.
For $k=2$, $E=(A,I_d-A)$ and $H=(B,-B)$, the restriction is exactly
\begin{equation}\label{eq:binarycovariance83}
 (\mathcal C_EH)_1=A(I_d-A)B+BA(I_d-A),\qquad
 \mathcal Q_N(E,F)^2=2Q_N(A,F_1)^2,
\end{equation}
where $Q_N$ is the binary comparison modulus in
\eqref{eq:matrixmodulus76}.
\end{theorem}

Thus the binary variance operator has a normalized-tuple extension,
not an independently regularized list of effects.  The theorem gives
a global upper certificate.  The two-sided balanced-interior
classification and its entropy theorem remain
Theorems~\ref{thm:multioutcomemetric82} and
\ref{thm:multioutcomeentropy82}; a matching general lower bound for
\eqref{eq:covariancemidpoint83} is a separate question.

\subsection{Positive covariance and horizontal decomposition}

\begin{lemma}[Covariance identities]
\label{lem:covarianceidentities83}
On all Hermitian tuples, $\mathcal C_E$ is self-adjoint and positive
semidefinite, and
\begin{align}
 \langle K,\mathcal C_EK\rangle
 &=\sum_j\tr(E_jK_j^2)-\|S_E(K)\|_{\rm HS}^2\notag\\
 &=\sum_j\|E_j^{1/2}(K_j-S_E(K))\|_{\rm HS}^2.
                                      \label{eq:covariancevariance83}
\end{align}
For $E_t=(1-t)E+tF$, $0\le t\le1$,
\begin{equation}\label{eq:covarianceconcavity83}
 \langle K,[\mathcal C_{E_t}-(1-t)\mathcal C_E-t\mathcal C_F]K\rangle
 =t(1-t)\|S_F(K)-S_E(K)\|_{\rm HS}^2.
\end{equation}
The kernel consists precisely of tuples satisfying
$E_j^{1/2}(K_j-S_E(K))=0$ for every $j$.
Moreover $\mathcal C_E=0$ on $\mathcal T_{d,k}$ if and only if the
$E_j$ are mutually orthogonal projections, with zero projections
permitted.
\end{lemma}
\begin{proof}
Write $A_j=E_j^{1/2}$ and let $A$ be the vertical stack of the $A_j$.
Then $A^*A=I_d$, so $P=I-AA^*$ is an orthogonal projection.  On
complex matrix tuples with real Hilbert--Schmidt inner product,
define
\[
 (T_AB)_j=A_j^*B_j+B_j^*A_j.
\]
Its real adjoint is $(T_A^*K)_j=2A_jK_j$.  Direct multiplication gives
\begin{equation}\label{eq:covariancefactor83}
 \mathcal C_E=\tfrac14 T_APT_A^*.
\end{equation}
This proves self-adjointness, positivity, and
\eqref{eq:covariancevariance83}.  It also proves the kernel claim.
Constant tuples are killed and the displayed formula has zero sum,
so no extra projection onto $\mathcal T_{d,k}$ is implicit.

In the first line of \eqref{eq:covariancevariance83}, the first term
is linear in $E$ and $S_E(K)$ is linear in $E$.  Expanding its squared
norm gives \eqref{eq:covarianceconcavity83}.

If the restriction of $\mathcal C_E$ to $\mathcal T_{d,k}$ vanishes,
subtracting the arithmetic mean from any tuple shows that its form
vanishes on all tuples.  Take $K_j=I_d$ at one index and zero
elsewhere.  Its form is $\tr(E_j-E_j^2)=0$.  Since $0\preceq E_j
\preceq I_d$, this forces $E_j^2=E_j$.  Their sum being $I_d$, they
are mutually orthogonal: each positive term in
$\sum_{i\ne j}E_jE_iE_j=0$ must vanish.  Conversely, for orthogonal
projections the summands $E_jK_j$ have mutually orthogonal ranges,
and $\|\sum_jE_jK_j\|_{\rm HS}^2=\sum_j\|E_jK_j\|_{\rm HS}^2$.
Their covariance form is zero.
\end{proof}

\begin{lemma}[Regularized horizontal energy]
\label{lem:horizontalenergy83}
For $H\in\mathcal T_{d,k}$ and $\tau>0$,
\begin{equation}\label{eq:energymin83}
 \langle H,(\mathcal C_E+\tau\operatorname{Id})^{-1}H\rangle
 =\min_{\substack{A^*B=0,\ R\in\mathcal T_{d,k}\\T_AB+R=H}}
       \left(4\|B\|_{\rm HS}^2+\tau^{-1}\|R\|_{\rm HS}^2\right).
\end{equation}
If $K=(\mathcal C_E+\tau\operatorname{Id})^{-1}H$, a minimizer is
\begin{equation}\label{eq:energyoptimizer83}
 B_j=\tfrac12 A_j(K_j-S_E(K)),\qquad R=\tau K.
\end{equation}
\end{lemma}
\begin{proof}
The constraint $A^*B=0$ means $PB=B$.  Thus
$B=PT_A^*K/4$ gives $T_AB=\mathcal C_EK$ and is feasible with
$R=\tau K$.  Its objective is
$\langle K,\mathcal C_EK\rangle+\tau\|K\|^2=\langle H,K\rangle$.
For any feasible perturbation $(\Delta B,\Delta R)$, its linear
objective term is
\[
 8\langle B,\Delta B\rangle+2\tau^{-1}\langle R,\Delta R\rangle
 =2\langle K,T_A\Delta B+\Delta R\rangle=0.
\]
The remaining terms are nonnegative.  This proves minimality,
including at singular $E$; no inverse of an effect was used.
\end{proof}

\subsection{Adaptive differentiation and boundary passage}

\begin{proof}[Proof of Theorem~\ref{thm:closedcovariance83}]
First suppose every effect at the point of differentiation is
positive definite.  The tangent $H_0=T_AB$ in
\eqref{eq:energyoptimizer83} has zero sum.  For small real $s$,
$E_j+s(H_0)_j$ is positive and the tuple remains normalized.
If $R_j(s)$ is its positive square root, then
\[
 U_j=(B_j-R_j'(0))A_j^{-1}
\]
is anti-Hermitian.  Indeed, multiplication of $U_j+U_j^*$ on the
left and right by $A_j$ gives
$A_jB_j+B_j^*A_j-A_jR_j'(0)-R_j'(0)A_j=0$.
The factors $\exp(sU_j)R_j(s)$ therefore realize the derivative
$B_j$ while preserving the effects exactly.

Use the measurement isometry with an observed label and an
inaccessible duplicate label, together with the factor environment.
At the point under consideration it obeys
\[
 W^*\dot W=\sum_jA_j^*B_j=0,\qquad
 \|\dot W\|_{\rm op}\le\|B\|_{\rm HS}.
\]
Purify any fixed common $N$-slot tester.  In the inner product of
terms differentiated at distinct slots, cancel the common later
isometries; the later differentiated slot contains $W^*\dot W$ or
its adjoint and vanishes.  Thus the observed output derivative in
the direction $H_0$ has trace norm at most
$2\sqrt N\|B\|_{\rm HS}$.  This holds for every reference and
memory dimension.  The environments used for comparison are not
returned to the tester.

The output derivative of the same tester is linear in the channel
tangent.  For the residual tuple $R$, one call on a joint state
$\rho$ has reference blocks $\tr_A[(R_j\otimes I)\rho]$.
Duality against Hermitian contractions bounds the trace norm of the
$j$th block by $\|R_j\|_{\rm op}$.  Inserting the residual at each
slot and using trace-norm contraction hence costs at most
\[
 N\sum_j\|R_j\|_{\rm op}\le N\sqrt{k}\|R\|_{\rm HS}.
\]
This is a differential decomposition, not a proposed physical mixture
of measurements.  With $\tau=N^{-1}$, Cauchy--Schwarz and
\eqref{eq:energymin83} now give
\begin{align*}
 2\sqrt N\|B\|_{\rm HS}+N\sqrt{k}\|R\|_{\rm HS}
 &\le\sqrt{k+1}
       \sqrt{4N\|B\|_{\rm HS}^2+N^2\|R\|_{\rm HS}^2}\\
 &=\sqrt{k+1}\,g_{N,E}(H).
\end{align*}
Integrate for the fixed tester, then take the supremum.  Padding a
publicly stopped experiment with discarded dummy calls makes its
original stopped output a common postprocessing of an $N$-slot
experiment, so the same bound applies.

For a boundary path, use $E^\epsilon(t)=(1-\epsilon)E(t)+\epsilon U$,
where $U_j=I_d/k$ and $0<\epsilon<1$.  The covariance is polynomial
and positive semidefinite, so its regularized inverse is continuous.
Moreover $g_{N,E}(H)\le N\|H\|_{\rm HS}$, uniformly on the closed
body.  Dominated convergence applies to the integrals.  The hybrid
bound $D_N(E,F)\le N\sum_j\|E_j-F_j\|_{\rm op}$ gives convergence
of the endpoint distances.  This proves \eqref{eq:covariancepath83}
on the entire body.

For the segment $E(t)=(1-t)E+tF$, put $M=(E+F)/2$ and $H=F-E$.
Concavity and positivity imply, for $0<t\le1/2$,
\[
 \mathcal C_{E(t)}+\tau\operatorname{Id}
 \succeq2t(\mathcal C_M+\tau\operatorname{Id}).
\]
Inverse order bounds $g_{N,E(t)}(H)$ by $(2t)^{-1/2}g_{N,M}(H)$;
use $2(1-t)$ on the second half.  The two scalar integrals sum to
$2$, proving \eqref{eq:closedcovariance83}.  Finally, substituting
$K=(B,-B)$ into \eqref{eq:covarianceoperator83} gives
\eqref{eq:binarycovariance83}.  The factor two in the quadratic form
comes from the two equal Hilbert--Schmidt contributions.
\end{proof}

\subsection{Classical postprocessing}

\begin{proposition}[Data processing of horizontal energy]
\label{prop:covarianceprocessing83}
Let $T=(T_{aj})$ be a column-stochastic classical channel and
$F_a=\sum_jT_{aj}E_j$.  Write $(T^*L)_j=\sum_aT_{aj}L_a$.
Then, on all Hermitian tuples,
\begin{equation}\label{eq:covarianceprocessing83}
 \langle L,\mathcal C_FL\rangle
 \ge\langle T^*L,\mathcal C_E(T^*L)\rangle.
\end{equation}
Consequently the extended quadratic energy
\[
 \mathcal I_E(H)=\sup_K
       \{2\langle H,K\rangle-\langle K,\mathcal C_EK\rangle\}
\]
satisfies $\mathcal I_F(TH)\le\mathcal I_E(H)$.  The value may be
infinite at a boundary; when finite it is
$\langle H,\mathcal C_E^\dagger H\rangle$.
\end{proposition}
\begin{proof}
The two $S$ matrices coincide.  The difference of the covariance
forms is therefore
\[
 \sum_j\tr\!\left[E_j\left(\sum_aT_{aj}L_a^2
                   -(\sum_aT_{aj}L_a)^2\right)\right]\ge0.
\]
The bracket equals the positive sum
$\tfrac12\sum_{a,b}T_{aj}T_{bj}(L_a-L_b)^2$.
Substitution in the variational definition and enlargement from
$K=T^*L$ to all $K$ prove the energy inequality.  The pseudoinverse
statement follows by diagonalizing the real self-adjoint covariance;
a component along its kernel makes the supremum infinite.
\end{proof}

\subsection{A uniform noise crossover, sharp in horizon and noise}

\begin{theorem}[Regularization of arbitrary support patterns]
\label{thm:noisecrossover83}
For $0\le\epsilon\le1$ put $E^\epsilon=(1-\epsilon)E+\epsilon U$.
On $\mathcal T_{d,k}$,
\begin{equation}\label{eq:covariancenoisefloor83}
 \mathcal C_{E^\epsilon}\succeq
       (1-\epsilon)\mathcal C_E+\frac\epsilon k\operatorname{Id}.
\end{equation}
In particular, for arbitrary $E,F\in\mathcal P_{d,k}$,
\begin{equation}\label{eq:noisecrossoverupper83}
 D_N(E^\epsilon,F^\epsilon)\le
 \min\left\{2,
 \frac{2\sqrt{k+1}(1-\epsilon)N}
      {\sqrt{1+N\epsilon/k}}\|F-E\|_{\rm HS}\right\}.
\end{equation}
The dependence on $N$ and $\epsilon$ is attained, up to constants
depending only on $k$, by a two-dimensional family.  Namely, let
$P=|0\rangle\langle0|$, $P_\theta=|u_\theta\rangle\langle u_\theta|$,
$u_\theta=\cos\theta|0\rangle+\sin\theta|1\rangle$, and
\[
 E=(P,I_2-P,0,\ldots,0),\qquad
 F=(P_\theta,I_2-P_\theta,0,\ldots,0).
\]
For $0\le\theta\le\pi/2$, $0\le\epsilon\le1/2$, and
$z=(1-\epsilon)N\sin\theta/\sqrt{1+N\epsilon}$,
\begin{equation}\label{eq:noisecrossoverlower83}
 \frac{\min\{1,z\}}{16384}
 \le D_N^{\rm na}(E^\epsilon,F^\epsilon)
 \le D_N(E^\epsilon,F^\epsilon)
 \le4\sqrt{k(k+1)}\min\{1,z\}.
\end{equation}
\end{theorem}
\begin{proof}
For $K\in\mathcal T_{d,k}$, $S_U(K)=0$ and
$\mathcal C_UK=K/k$.  Concavity proves
\eqref{eq:covariancenoisefloor83}.  Apply it at the midpoint and use
Theorem~\ref{thm:closedcovariance83} to obtain
\eqref{eq:noisecrossoverupper83}.

For the displayed family, keep labels one and two and send every
other label to an independent fair bit.  This common postprocessing
produces the binary effects
$B=\beta P+\epsilon I_2/2$ and
$B_\theta=\beta P_\theta+\epsilon I_2/2$, where $\beta=1-\epsilon$.
The equality holds on subnormalized reference states.  It therefore
preserves the nonadaptive lower witnesses as well as channel
postprocessing.  The midpoint eigenvalues are
$(1\pm\beta\cos\theta)/2$, and direct substitution in
\eqref{eq:matrixcoordinates76} gives
\[
 Q_N(B,B_\theta)^2
 =\frac{4N^2\beta^2\sin^2\theta}
             {2+N(1-\beta^2\cos^2\theta)}
 \ge\frac{4z^2}{2+z^2}.
\]
Here $1-\beta^2\le2\epsilon$ and $N\ge1$ were used.
It follows that $\min\{1,Q_N(B,B_\theta)\}\ge\min\{1,z\}$.
The binary lower theorem in dimension two gives the first inequality
of \eqref{eq:noisecrossoverlower83}.  For the upper,
$\|F-E\|_{\rm HS}=2\sin\theta$ and
$(1+N\epsilon)/(1+N\epsilon/k)\le k$.
Combine \eqref{eq:noisecrossoverupper83} with $D_N\le2$.
\end{proof}

The family passes from projections and zero effects at $\epsilon=0$
to strictly positive effects at every $\epsilon>0$.  Thus a bound
uniform over the closed body must accommodate both the $N$ scale
and the $\sqrt{N/\epsilon}$ scale.  This conclusion does not assert a
matching entropy law on all multi-outcome boundary strata.

\subsection{Exact evaluation without effect square roots}

\begin{proposition}[Polynomial-size rational certificate]
\label{prop:covarianceexact83}
For Gaussian-rational $E,F\in\mathcal P_{d,k}$ and integer $N\ge1$,
$\mathcal Q_N(E,F)^2$ is rational and can be evaluated by exact
linear algebra on a positive definite real system of dimension
$p=(k-1)d^2$.  The bit complexity is polynomial in $d,k$, the total
input bit length, and $\log N$.  Legality is decidable by exact
positive-semidefinite tests.  No effect square root, ODE integration,
rank decision by tolerance, or enumeration of a dictionary is required.
\end{proposition}
\begin{proof}
Use the rational Hermitian matrix basis consisting of diagonal
units and real and imaginary off-diagonal pairs.  For each of the
first $k-1$ slots, insert one matrix basis element there and its
negative in the last slot.  These tuples $B_a$ form a rational basis
of $\mathcal T_{d,k}$; they need not be orthonormal.  Form
\[
 G_{ab}=\langle B_a,B_b\rangle,\quad
 L_{ab}=\langle B_a,\mathcal C_{(E+F)/2}B_b\rangle,\quad
 b_a=\langle B_a,F-E\rangle.
\]
All entries are rational, $G$ is positive definite, and $L\succeq0$.
Solve $(L+N^{-1}G)x=b$.  Then
$\mathcal Q_N^2=N b^{\mathsf T}x$.  Omitting the Gram matrix would
compute a different regularization.  Fraction-free elimination
performs polynomially many rational arithmetic operations; clearing
input denominators and applying determinant bounds gives polynomial
bit length for all minors and for the output.  The same reasoning
applies to exact symmetric elimination for each effect's positivity.
The complexity statement concerns this certificate evaluator, not
the exhaustive public code or collective-readout synthesis.
\end{proof}

\section{A finite product estimate}\label{sec:product85}
All classical distances in this paper are unhalved $\ell^1$ distances.
The following elementary estimate is the Bernoulli-product input to
Theorems~\ref{thm:orbittrichotomy84} and \ref{thm:multioutcomemetric82}.
We include its proof so that the support and discrimination theorems
can be checked without the binary supplement.

\begin{lemma}[A uniform finite Bernoulli lower bound]\label{lem:bernoulli85}
For $p,q\in[0,1]$ and every integer $N\ge1$,
\begin{equation}\label{eq:bernoulli85}
 \|\operatorname{Ber}(p)^{\otimes N}-
       \operatorname{Ber}(q)^{\otimes N}\|_1
 \ge \frac1{128}\min\{1,\sqrt N|p-q|\}.
\end{equation}
In particular, if a differentiable two-sided probability curve satisfies
$p'(0)\ne0$, then, on a sufficiently small two-sided interval,
its product experiment separates $p(0)$ from $p(t)$ by at least
$c\min\{1,\sqrt N|t|\}$ for every $N$, with $c>0$ depending
on that curve.  Necessarily $0<p(0)<1$.
\end{lemma}
\begin{proof}
Set $h=|p-q|$.  First let $m\ge1$ and $\sqrt m h\le1/4$.
For $u=m^{-1/2}$ define $z(a)=1-a+ae^{iu}$, $0\le a\le1$.
The argument of $z(a)$ has a continuous choice since its real part
is positive.  Direct differentiation and $0<u\le1$ give
\[
 |z(a)|^2\ge1-\frac1{4m},\qquad
 \frac u2\le\frac{d}{da}\arg z(a)
       =\frac{\sin u}{|z(a)|^2}\le\frac{4u}3.
\]
Bernoulli's inequality gives $|z(a)|^m\ge\sqrt{3/4}>3/4$.
The arguments of $z(p)^m$ and $z(q)^m$ differ by an angle
$\vartheta\in[\sqrt m h/2,4\sqrt m h/3]\subset[0,1/3]$.
For two complex numbers of moduli at least $3/4$,
\[
 |z(p)^m-z(q)^m|\ge\tfrac32\sin(\vartheta/2)
 \ge\frac{3}{2\pi}\vartheta\ge\frac1{16}\sqrt m h.
\]
These complex numbers are the expectations of
$\exp(iu\sum_{a=1}^m X_a)$ under the two Bernoulli product laws.
This statistic has modulus one, so their $\ell^1$ distance is at least
its expectation difference.

If $h\le1/4$ and $\sqrt N h\le1/4$, use $m=N$.
If $h\le1/4$ and $\sqrt N h>1/4$, use
$m=\lfloor(16h^2)^{-1}\rfloor\le N$.
Then $1/(4\sqrt2)\le\sqrt m h\le1/4$, since
$\lfloor x\rfloor\ge x/2$ for $x\ge1$.
Discarding the other observations is a common processing and gives
a lower bound $1/(64\sqrt2)>1/128$.
For $h>1/4$ the one-copy distance is $2h>1/2$.
The case $h=0$ is immediate.  Finally differentiability gives
$|p(t)-p(0)|\ge|p'(0)||t|/2$ near zero.  A two-sided probability
curve cannot have a nonzero first derivative at either endpoint.
\end{proof}

\section{Support geometry of the covariance kernel}\label{sec:support84}
Let $\mathcal H_d$ be the real Hilbert space of Hermitian $d\times d$
matrices, with inner product $\tr(AB)$.  Throughout this section,
$E=(E_1,\ldots,E_k)$ is an arbitrary ordered measurement:
$E_j\succeq0$ and $\sum_jE_j=I_d$.  Put
\[
 P_j=\operatorname{supp}E_j,\qquad Q_j=I_d-P_j,\qquad r_j=\operatorname{rank}E_j.
\]
The support of a zero effect is the zero projection.  The real subspaces
\begin{equation}\label{eq:supportspan84}
 \mathcal W_E=\sum_{j=1}^k P_j\mathcal H_dP_j,
 \qquad \mathcal A_E=\mathcal W_E^\perp
   =\{A\in\mathcal H_d:P_jAP_j=0\text{ for every }j\}
\end{equation}
will be called the \emph{support span} and its orthogonal complement.
The sum in \eqref{eq:supportspan84} is a sum of real linear subspaces,
not an assertion that $\mathcal W_E$ is an algebra.  In particular,
for higher-rank effects it need not equal the span of the effects themselves.
We use the zero-sum tuple space
$\mathcal T=\{K\in\mathcal H_d^k:\sum_jK_j=0\}$ with the sum
Hilbert--Schmidt inner product, and the covariance of
Lemma~\ref{lem:covarianceidentities83}.

\begin{theorem}[Complete support description of the kernel]\label{thm:supportkernel84}
For $A\in\mathcal A_E$ and $Z_j\in Q_j\mathcal H_dQ_j$, set
\begin{equation}\label{eq:kernelparam84}
 X_j=i[P_j,A]+Z_j,\qquad
 K_j=X_j-\frac1k\sum_{\ell=1}^kX_\ell.
\end{equation}
This is a real linear bijection from
$\mathcal A_E\oplus\bigoplus_jQ_j\mathcal H_dQ_j$ onto
$\ker(\mathcal C_E|_{\mathcal T})$.  Consequently
\begin{equation}\label{eq:kerneldimension84}
 \dim\ker(\mathcal C_E|_{\mathcal T})
 =d^2-\dim\mathcal W_E+\sum_j(d-r_j)^2.
\end{equation}
In particular the kernel, and not just its dimension, depends only
on the ordered support projections.
\end{theorem}
\begin{proof}
Write $S=\sum_jE_jK_j$.  The nonnegative variance decomposition in
Lemma~\ref{lem:covarianceidentities83} shows that $K$ is in the kernel
exactly when
\begin{equation}\label{eq:supportkernelcondition84}
 P_jK_j=P_jS\quad\text{for every }j.
\end{equation}
Indeed multiplication by $\sqrt{E_j}$ has precisely the same kernel
as multiplication by $P_j$; no inverse at a singular effect is used.
Decompose $S=B+iA$, where $A,B$ are Hermitian.  Since
$P_jK_jP_j$ and $P_jBP_j$ are Hermitian, compressing
\eqref{eq:supportkernelcondition84} gives $P_jAP_j=0$.
Furthermore $P_j(K_j-B)=iP_jA$, with the adjoint identity on the
other side.  These determine the two off-diagonal support blocks
and the support block of $K_j-B$; its remaining block is arbitrary.
Thus
\[
 K_j=B+i[P_j,A]+Z_j,\qquad Z_j=Q_jZ_jQ_j=Z_j^*.
\]
The condition $\sum_jK_j=0$ forces
$B=-k^{-1}\sum_j(i[P_j,A]+Z_j)$, proving necessity.

Conversely, let $A,Z_j$ be as stated.  Because $P_jAP_j=0$,
$E_jAP_j=0$, and hence
\[
 \sum_jE_jX_j=i\sum_jE_jA=iA.
\]
For the mean-subtracted tuple, therefore, $S=iA-\overline X$,
where $\overline X=k^{-1}\sum_jX_j$ is Hermitian.
We have $P_jX_j=iP_jA$, which proves
\eqref{eq:supportkernelcondition84}.  To check injectivity, suppose
that all $K_j$ vanish.  Then all $X_j$ equal a common Hermitian
matrix $D$, so the last display gives $D=iA$.  A matrix that is
both Hermitian and anti-Hermitian is zero.  Therefore $A=D=0$ and
all $Z_j=0$.  Finally $\dim(Q_j\mathcal H_dQ_j)=(d-r_j)^2$;
rank--nullity in \eqref{eq:supportspan84} proves
\eqref{eq:kerneldimension84}.
\end{proof}

\begin{proposition}[Unitary directions and the support span]\label{prop:unitaryrange84}
Let $G\in\mathcal H_d$ and $H_j=i[G,E_j]$.  Then
\begin{equation}\label{eq:unitaryrange84}
 H\in\operatorname{Ran}(\mathcal C_E|_{\mathcal T})
 \quad\Longleftrightarrow\quad G\in\mathcal W_E.
\end{equation}
More precisely, the pairing of $H$ with the kernel tuple
\eqref{eq:kernelparam84} is $-2\tr(GA)$.
\end{proposition}
\begin{proof}
We have $Q_jH_jQ_j=0$, so the $Z_j$ terms have zero pairing with $H$;
its zero tuple sum also removes $\overline X$.  Trace cyclicity gives
\[
 \sum_j\tr\bigl(H_ji[P_j,A]\bigr)
 =\tr\left(A\sum_ji[H_j,P_j]\right).
\]
Expanding in the displayed order yields
\[
 \sum_ji[H_j,P_j]
 =-2G+\sum_j(E_jGP_j+P_jGE_j).
\]
Every summand in the last sum belongs to
$P_j\mathcal H_dP_j$; it is therefore orthogonal to $A$.
The pairing is $-2\tr(GA)$ as asserted.  A self-adjoint linear
operator in finite dimension has range equal to the orthogonal
complement of its kernel.  Theorem~\ref{thm:supportkernel84} now
proves \eqref{eq:unitaryrange84}.
\end{proof}

\begin{corollary}[Regularized tangent scales]\label{cor:tangentscales84}
Let $\Pi_0$ be the orthogonal projection onto
$\ker(\mathcal C_E|_{\mathcal T})$, and let $u_\ell$ be an orthonormal
basis of its complement with $\mathcal C_Eu_\ell=\lambda_\ell u_\ell$,
$\lambda_\ell>0$.  For every $H\in\mathcal T$,
\begin{equation}\label{eq:tangentspectral84}
 g_{N,E}(H)^2
 =N^2\|\Pi_0H\|_{\mathrm{HS}}^2+
  N\sum_\ell\frac{|\langle H,u_\ell\rangle|^2}{\lambda_\ell+N^{-1}}.
\end{equation}
Thus a nonzero horizontal direction has order $\sqrt N$, whereas a
direction with nonzero kernel component has order $N$, at a fixed
measurement.  For $H=i[G,E]$ the first alternative is exactly
$G\in\mathcal W_E$ and $H\ne0$.
\end{corollary}
\begin{proof}
Apply the finite-dimensional spectral theorem to
$N(\mathcal C_E+N^{-1}\operatorname{Id})^{-1}$ and use
Proposition~\ref{prop:unitaryrange84}.
\end{proof}
The corollary concerns a comparison form, not yet an operational
lower bound.  The distinction matters at the boundary; the next
section proves the corresponding finite-pair statement on entire
unitary orbits.

\begin{proposition}[Exact represented-input decision]\label{prop:supportexact84}
On Gaussian-rational effects and a Gaussian-rational Hermitian
generator, membership in $\mathcal W_E$, stationarity of the orbit,
and the kernel dimension \eqref{eq:kerneldimension84} are decidable
by exact rational linear algebra in polynomial bit complexity.
The same calculation gives the orthogonal projection of $G$ onto
$\mathcal W_E$.
\end{proposition}
\begin{proof}
First use exact positive-semidefinite and normalization tests to
validate the input.  Select independent columns $V_j$ of $E_j$ by
rational elimination.  For a nonzero effect its support is
\[
 P_j=V_j(V_j^*V_j)^{-1}V_j^*;
\]
for a zero effect use $P_j=0$.  Compress a rational Hermitian basis
by all $P_j$ and select a real-linearly independent subfamily
$W_1,\ldots,W_q$.  The positive Gram system
$\sum_b\tr(W_aW_b)c_b=\tr(W_aG)$ gives
$\Pi_{\mathcal W_E}G=\sum_bc_bW_b$.  This tests membership exactly.
Commutators test stationarity.  Formula~\eqref{eq:kerneldimension84}
then gives the dimension.  There are at most $kd^2$ candidate
matrices and $d^2$ independent matrices.  Clearing denominators
and applying fraction-free elimination gives polynomial bit
bounds by the determinant bounds used in
Proposition~\ref{prop:covarianceexact83}.  This is a decision about
supports and a specified tangent; no optimum over testers or
quantum-control synthesis is computed.
\end{proof}

\section{Finite-angle distinguishability on unitary orbits}\label{sec:orbit84}
Fix an ordered measurement $E$ and a Hermitian generator $G$, and write
\begin{equation}\label{eq:unitorbit84}
 E_j(t)=e^{itG}E_je^{-itG},\qquad H_j=i[G,E_j].
\end{equation}
Both $E$ and $G$ are known in this pair-comparison problem.  A tester
must use the same controls under the two hypotheses $E$ and $E(t)$;
its design may depend on this known pair.  It has access to the
classical device labels and its own reference and memory, never to
an unobserved measurement environment.

The Hamiltonian-not-in-Kraus-span criterion and its quantum-error-
correction interpretation are established in quantum channel metrology
\cite[Theorems~1--3]{ZJ84}.  The support criterion below is its
measurement-channel specialization.  We prove a finite-angle,
finite-use trace-distance statement directly, including an explicit
accumulated error estimate.  Thus no lower bound is inferred solely
from an infinitesimal Fisher-information scaling.

\begin{theorem}[Finite-use trichotomy on every fixed unitary orbit]\label{thm:orbittrichotomy84}
Let $d,N\ge1$, $k\ge2$, and let $E$ and $G$ be as in
\eqref{eq:unitorbit84}, with no rank or strict-positivity assumption.
The following alternatives are exhaustive.
\begin{enumerate}
\item If $[G,E_j]=0$ for every $j$, then $D_N(E,E(t))=0$ for all
$N$ and all $t$.
\item If $H\ne0$ and $G\in\mathcal W_E$, there are constants
$c,C,t_0>0$, depending on $E,G$, such that
\begin{equation}\label{eq:orbitsql84}
 c\min\{1,\sqrt N|t|\}\le D_N(E,E(t))
 \le C\min\{1,\sqrt N|t|\},\qquad |t|\le t_0.
\end{equation}
The lower bound is attained by a product-input, nonadaptive experiment.
\item If $G\notin\mathcal W_E$, there are constants
$c,C,t_0>0$, depending on $E,G$, such that
\begin{equation}\label{eq:orbithl84}
 c\min\{1,N|t|\}\le D_N(E,E(t))
 \le C\min\{1,N|t|\},\qquad |t|\le t_0.
\end{equation}
The lower bound admits a common adaptive experiment carrying a logical
qubit between calls, with a reference of dimension at most $2d$ at
an individual call and ideal trusted controls.
\end{enumerate}
The constants are not uniform as the measurement or generator changes.
All upper bounds allow arbitrary retained reference and adaptive memory,
and bounded public stopping.  This theorem classifies a fixed unitary
orbit; it does not assert two-sided equivalence of the midpoint
covariance modulus on arbitrary pairs in the complete measurement body.
\end{theorem}

We first record the error-correction statement in the form needed for
finite angles.  Its exact-correction premise is the usual
Knill--Laflamme condition \cite{KL84}; a proof is included to specify
which systems are available to the recovery.

\begin{lemma}[A corrected logical channel and its quadratic error]\label{lem:correctedchannel84}
Let $\mathcal M$ be a channel on an input space $\mathsf A$, and
let $J:\C^2\longrightarrow\mathsf A\otimes\mathsf R$ be an isometry.
Suppose a channel $\mathcal R$ from the actual output of
$\mathcal M\otimes\operatorname{Id}_{\mathsf R}$ to $\C^2$ satisfies
\[
 \mathcal R\circ(\mathcal M\otimes\operatorname{Id}_{\mathsf R})
          \circ\mathcal J=\operatorname{Id}_2,
 \qquad \mathcal J(X)=JXJ^*.
\]
For a Hermitian $G$ on $\mathsf A$, put
$G_L=J^*(G\otimes I)J$, $\mathcal U_t(X)=e^{-itG_L}Xe^{itG_L}$,
and
\[
 \Phi_t=\mathcal R\circ(\mathcal M\otimes\operatorname{Id}_{\mathsf R})
       \circ\operatorname{Ad}(e^{-it(G\otimes I)})\circ\mathcal J.
\]
In unhalved diamond norm,
\begin{equation}\label{eq:logicalerror84}
 \|\Phi_t-\mathcal U_t\|_\diamond
 \le4t^2\|G\|_{\mathrm{op}}^2,
 \qquad
 \|\Phi_t^{\circ m}-\mathcal U_t^{\circ m}\|_\diamond
 \le4mt^2\|G\|_{\mathrm{op}}^2.
\end{equation}
\end{lemma}
\begin{proof}
Let $\mathcal T=\mathcal R\circ(\mathcal M\otimes\operatorname{Id})$
and choose Kraus operators $V_\ell$ for $\mathcal T$.  Since
$\mathcal T\circ\mathcal J$ is the identity channel, whose Choi
operator has rank one, $V_\ell J=c_\ell I_2$ with
$\sum_\ell|c_\ell|^2=1$.  Trace preservation on the entire input
space gives
\[
 \sum_\ell c_\ell^*V_\ell
 =\sum_\ell(V_\ell J)^*V_\ell
 =J^*\sum_\ell V_\ell^*V_\ell=J^*.
\]
It follows by differentiating that $\Phi_0=\operatorname{Id}_2$
and $\Phi'_0=-i[G_L,\,\cdot\,]$.
For any Hermitian $B$, the generator of $\operatorname{Ad}(e^{-itB})$
has diamond norm at most $2\|B\|_{\mathrm{op}}$.  Its second
 derivative has diamond norm at most $4\|B\|_{\mathrm{op}}^2$,
since unitary channels have diamond norm one.  Taylor's integral
remainder and contractivity under the surrounding channels bound
the remainder for $\Phi_t$ by $2t^2\|G\|_{\mathrm{op}}^2$.
The remainder for $\mathcal U_t$ has the same bound because
$\|G_L\|_{\mathrm{op}}\le\|G\|_{\mathrm{op}}$.  Their zeroth
and first derivatives agree, proving the first inequality.
Telescoping the two $m$-fold compositions of channels proves the second.
\end{proof}

\begin{lemma}[A reference-preserving code from the support complement]\label{lem:supportcode84}
Suppose $G\notin\mathcal W_E$ and put
\begin{equation}\label{eq:signstates84}
 A=G-\Pi_{\mathcal W_E}G,\qquad
 s=\tr A_+=\tr A_->0,\qquad \rho_\pm=A_\pm/s.
\end{equation}
There is a code $J:\C^2\to\C^d\otimes\mathsf R$ with
$\dim\mathsf R\le2d$, corrected exactly by a channel acting only
on the classical label and $\mathsf R$ after a call to $E$.
The logical generator $G_L=J^*(G\otimes I)J$ is diagonal with gap
\begin{equation}\label{eq:logicalgap84}
 \Delta=\tr\bigl(G(\rho_+-\rho_-)\bigr)
        =\|A\|_{\mathrm{HS}}^2/s>0.
\end{equation}
\end{lemma}
\begin{proof}
The identity belongs to $\mathcal W_E$ because $\sum_jE_j=I$;
thus $A$ is traceless.  Its nonzero positive and negative parts have
equal trace.  Also
\[
 P_j\rho_+P_j=P_j\rho_-P_j\quad\text{for every }j.
\]
Purify $\rho_+$ and $\rho_-$ in $\C^d\otimes\C^d$ and attach
orthogonal flags to the second factor.  Map the two logical basis
vectors to these flagged purifications.  They are orthonormal,
and the flags make all off-diagonal code matrix elements of
$B\otimes I$ vanish.  The two diagonal elements are equal for
every $B\in P_jM_d(\C)P_j$ by the compression identity above.

Choose measurement Kraus operators
$L_{j,a}=|j\rangle\langle a|\sqrt{E_j}$, $1\le a\le d$.
Products with different labels are zero; for one label their span
is $P_jM_d(\C)P_j$.  Hence, with $\mu=(j,a)$,
\begin{equation}\label{eq:kl84}
 J^*(L_\mu^*L_\nu\otimes I)J=\lambda_{\mu\nu}I_2
\end{equation}
for a positive matrix $\lambda$ of trace one.  Diagonalize
$\lambda$ by a unitary change of Kraus representation.  If its
positive eigenvalues are $\lambda_b$, the corresponding operators
$(\widetilde L_b\otimes I)J/\sqrt{\lambda_b}$ are isometries
with mutually orthogonal ranges.  A channel projects onto those
ranges and applies their inverse isometries; on the orthogonal
complement it prepares a fixed logical state.  This is completely
positive and trace preserving on the entire output space, and it
corrects the code exactly.  Its input space is just the actual
classical-label register tensor $\mathsf R$.  It does not contain
the lost input or a dilation environment.  Equivalently, since the
actual state is label diagonal, the same channel may first read the
label and then act on the retained reference.

The logical generator is diagonal by the flags, with entries
$\tr(G\rho_\pm)$.  Orthogonality of $A$ to
$\Pi_{\mathcal W_E}G$ gives
$\tr(GA)=\tr(A^2)$, proving \eqref{eq:logicalgap84}.
\end{proof}

\begin{proof}[Proof of Theorem~\ref{thm:orbittrichotomy84}]
The first assertion follows directly from \eqref{eq:unitorbit84}.
For the second, Proposition~\ref{prop:unitaryrange84} gives
\[
 h^2=\langle H,\mathcal C_E^\dagger H\rangle<\infty.
\]
Conjugating all tuple components by $e^{itG}$ intertwines the
covariances and their tangent directions, so this value is constant
along the orbit.  Inverse spectral comparison gives
$g_{N,E(t)}(E'(t))\le\sqrt N h$.  The complete-body path theorem
therefore yields
\begin{equation}\label{eq:orbitupper84}
 D_N(E,E(t))\le\min\{2,\sqrt{k+1}\,h\sqrt N|t|\}.
\end{equation}
There is an index $j$ with $H_j\ne0$.  Choose a unit vector $v$
with $\langle v,H_jv\rangle\ne0$.  The two Bernoulli parameters
obtained by input $v$ and the event ``label $j$'' have gap at least
$c_0|t|$ on a sufficiently small two-sided interval.  The finite
Bernoulli-product bound used in
Theorem~\ref{thm:multioutcomemetric82} gives
$(1/128)\min\{1,c_0\sqrt N|t|\}$.  This proves
\eqref{eq:orbitsql84}, including all finite $N$.

For the third alternative, use Lemma~\ref{lem:supportcode84}.
The equality
$\mathcal M_{E(t)}=\mathcal M_E\circ\operatorname{Ad}(e^{-itG})$
holds also with a retained reference.  Encode, make one device call,
recover, and repeat $m\le N$ times.  Under $E$ the logical channel
is the identity.  Under $E(t)$ it is $\Phi_t^{\circ m}$ from
Lemma~\ref{lem:correctedchannel84}.  Start in the equal superposition
of the two logical basis states.  The ideal logical unitary produces
unhalved trace distance $2|\sin(mt\Delta/2)|$ from that state.
Consequently, for every real $t$,
\begin{equation}\label{eq:finiteanglelower84}
 D_N(E,E(t))\ge
 \max_{1\le m\le N}
 \left(2|\sin(mt\Delta/2)|-4mt^2\|G\|_{\mathrm{op}}^2\right)_+.
\end{equation}
This bound concerns one specified common tester for each $m$,
not an inaccessible choice of a different recovery under the two hypotheses.

For $t\ne0$ choose
\[
 |t|\le t_0:=\min\left\{\frac1{2\Delta},
                       \frac{\Delta}{4\pi\|G\|_{\mathrm{op}}^2}\right\},
 \qquad
 m=\min\left\{N,\left\lfloor\frac1{|t|\Delta}\right\rfloor\right\}.
\]
Then $x=m|t|\Delta\le1$ and
$x\ge\tfrac12\min\{1,N|t|\Delta\}$.
Since $2\sin(x/2)\ge2x/\pi$ and the error in
\eqref{eq:finiteanglelower84} is at most $x/\pi$, we obtain
\begin{equation}\label{eq:explicitlocalhl84}
 D_N(E,E(t))\ge\frac1{2\pi}\min\{1,N|t|\Delta\}.
\end{equation}
The elementary hybrid estimate for the input unitaries gives
$D_N(E,E(t))\le\min\{2,2N|t|\|G\|_{\mathrm{op}}\}$,
which proves \eqref{eq:orbithl84}.  The zero-angle case is immediate.
The protocol uses deterministic $m\le N$ calls, hence belongs to
the class allowing bounded public stopping.  The upper estimates
hold for that entire class.  Finally if the first alternative held
with $G\notin\mathcal W_E$, then $H=0$ would contradict
Proposition~\ref{prop:unitaryrange84}; the alternatives are disjoint.
\end{proof}

\begin{corollary}[Full-rank, rank-one and projective measurements]\label{cor:orbitexamples84}
If one effect is positive definite, every nonconstant unitary orbit
has the square-root law \eqref{eq:orbitsql84}.  If all nonzero
effects have rank one, every nonconstant orbit has that law precisely
when the measurement is informationally complete, meaning that its
effects span $\mathcal H_d$.  If such a rank-one measurement is not
informationally complete, some generator has the linear law
\eqref{eq:orbithl84}.  For a projection-valued measurement every
orbit is either constant or has the linear law.
\end{corollary}
\begin{proof}
A full-rank support contributes $\mathcal H_d$ to $\mathcal W_E$.
For rank-one supports, $P_j\mathcal H_dP_j=\R E_j$, so that
$\mathcal W_E$ is exactly the effect span.  If this span is proper,
a nonzero element of its orthogonal complement is a generator
outside $\mathcal W_E$ and cannot commute with all effects by
Proposition~\ref{prop:unitaryrange84}.  For a projection-valued
measurement $\mathcal W_E$ is the block-diagonal Hermitian space
of the orthogonal supports, which is exactly the space of generators
commuting with every effect.  Apply Theorem~\ref{thm:orbittrichotomy84}.
\end{proof}

\begin{proposition}[Support spans under output processing]\label{prop:supportprocessing84}
Let $T_{aj}\ge0$ satisfy $\sum_aT_{aj}=1$, and put
$F_a=\sum_jT_{aj}E_j$.  Then $\mathcal W_E\subseteq\mathcal W_F$.
Thus classical processing of the ordered output cannot turn a
nonconstant square-root orbit into a linear orbit.  A reversible
splitting of labels preserves the support span and the classification.
\end{proposition}
\begin{proof}
For positive matrices, the kernel of a positive weighted sum is the
intersection of the kernels of its positively weighted summands.
Thus the support of $F_a$ contains $P_j$ whenever $T_{aj}>0$.
Every column has such a row, giving
$P_j\mathcal H_dP_j\subseteq\mathcal W_F$.
This proves the inclusion.  The orbit commutes with classical
processing, so Theorem~\ref{thm:orbittrichotomy84} applies with the
same generator.  A reverse stochastic map gives the opposite
inclusion for a reversible splitting.
\end{proof}

\begin{remark}[What the lower construction does not assert]\label{rem:orbitresources84}
The code and recovery depend on the known pair through $E,G$.
They do not constitute a common estimator of an unknown measurement.
Their existence is proved in finite-dimensional quantum mechanics;
polynomial gate synthesis, physical execution, and parameter-uniform
bounds for approximate controls are not inferred.  Formula
\eqref{eq:finiteanglelower84} quantifies the finite-angle channel
remainder under ideal controls.  It is not a claim of perfect
finite-use discrimination.  The full-body entropy and growing-$k$
learning questions are separate from this orbit classification.
\end{remark}

\section{Finite-use classification of smooth boundary curves}
\label{sec:curves85}
The support criterion extends beyond unitary conjugation.  Here
$E(t)=(E_1(t),\ldots,E_k(t))$, $|t|<a$, is a fixed two-sided
$C^2$ curve of ordered measurements, with $E=E(0)$ and
$H=E'(0)\ne0$.  The curve may change ranks at zero and need not
have a differentiable exact Kraus representation there.  All constants
below belong to this fixed curve, not uniformly to the measurement body.

Fix $0<a_0<a$.  Whenever a remainder below uses $L$, it means
\begin{equation}\label{eq:curvature86}
 L=\sup_{|u|\le a_0}\sum_{j=1}^k\|E_j''(u)\|_{\mathrm{op}}<\infty.
\end{equation}
This is the supremum of the sum, on the stated fixed interval.

\begin{lemma}[The support generator of a two-sided tangent]
\label{lem:curvegenerator85}
Let $P_j=\operatorname{supp}E_j$ and $Q_j=I-P_j$.  Then
$Q_jH_jQ_j=0$.  Conversely every Hermitian zero-sum tuple with these
vanishing blocks is the derivative of a real-analytic two-sided
measurement curve.  Define the Hermitian matrix
\begin{equation}\label{eq:curvegenerator85}
 \Gamma_E(H)=\frac i2\sum_j[P_j,H_j].
\end{equation}
On the real zero-sum tuple space,
\begin{equation}\label{eq:curverange85}
 H\in\operatorname{Ran}\mathcal C_E
 \quad\Longleftrightarrow\quad \Gamma_E(H)\in\mathcal W_E.
\end{equation}
For a unitary tangent $H_j=i[G,E_j]$,
$\Gamma_E(H)-G\in\mathcal W_E$.  In particular every stationary
unitary generator belongs to $\mathcal W_E$.
\end{lemma}
\begin{proof}
For $v\in\operatorname{Ran}Q_j$, the nonnegative function
$\langle v,E_j(t)v\rangle$ has a two-sided minimum at zero;
its derivative vanishes.  Polarization on that subspace gives
$Q_jH_jQ_j=0$.  The sum of the $H_j$ vanishes by normalization.
Pairing with the complete kernel tuple in
Theorem~\ref{thm:supportkernel84}, both the mean and the
missing-support blocks disappear, leaving
\[
 \langle H,K(A,Z)\rangle
 =\tr\left(A\sum_ji[H_j,P_j]\right)
 =-2\tr(A\Gamma_E(H)).
\]
Finite-dimensional self-adjointness proves \eqref{eq:curverange85}.
For a unitary tangent, expansion gives
\[
 \Gamma_E(H)=G-\tfrac12\sum_j(E_jGP_j+P_jGE_j),
\]
and the last sum belongs to $\mathcal W_E$.  If the tangent is zero,
this identity also proves the final assertion.  For the converse
realizability statement, the normalized-factor construction in the
proof of Lemma~\ref{lem:curvelogical85} is real analytic near zero
and has precisely the prescribed value and derivative.
\end{proof}

\begin{theorem}[Two-sided smooth-curve classification]\label{thm:curveclassification85}
For the fixed $C^2$ curve above there are $c,C,t_0>0$ such that,
for all integers $N\ge1$ and $|t|\le t_0$,
\begin{equation}\label{eq:curvedichotomy85}
 c\min\{1,N^\alpha|t|\}\le D_N(E,E(t))
       \le C\min\{1,N^\alpha|t|\},\qquad
 \alpha=
 \begin{cases}
  1/2,&\Gamma_E(H)\in\mathcal W_E,\\
  1,&\Gamma_E(H)\notin\mathcal W_E.
 \end{cases}
\end{equation}
The first lower bound uses repeated product inputs.  The second has a
known-pair adaptive realization with ideal encoding and recovery,
a logical qubit, and a per-call external reference of dimension at
most $2d$.  The same controls are used under both hypotheses.
Constants may degenerate with the tangent, nonzero support eigenvalues,
curvature, or support-complement gap.
\end{theorem}

The exponents and the error-correction mechanism have the established
channel-metrology antecedents \cite{ZJ84,KL84}.  The content here is a
direct finite-pair theorem from $C^2$ effect coordinates: the support
formula decides its branch even at second-order rank openings, without
assuming a smooth exact Kraus family or inferring a finite-angle lower
from Fisher information.  It is a theorem for each fixed nonzero
first-order curve, not a uniform arbitrary-pair midpoint equivalence.

\begin{lemma}[A horizontal curve with prescribed first derivative]
\label{lem:horizontalcurve85}
If $H\in\operatorname{Ran}\mathcal C_E$, there are factors
$A_j=\sqrt{E_j}$ and $B_j$ satisfying
\[
 A^*A=I,\quad A^*B=0,\quad
 A_j^*B_j+B_j^*A_j=H_j,
\]
where $A,B$ denote the vertically stacked matrices.  Put
$b=\|B\|_{\mathrm{op}}$ and $S=B^*B$.  The factors
\begin{equation}\label{eq:horizontalcurve85}
 A(t)=(A+tB)(I+t^2S)^{-1/2}
\end{equation}
define a normalized measurement $\widetilde E_j(t)=A_j(t)^*A_j(t)$
with the same zeroth and first derivatives as $E(t)$, and
\[
 A(t)^*A'(t)=0,\qquad
 \|A'(t)\|_{\mathrm{op}}\le b,\qquad
 \|A''(t)\|_{\mathrm{op}}\le7b^2\quad(|t|b\le1).
\]
If $L=\sup_{|s|\le a_0}\sum_j\|E_j''(s)\|_{\mathrm{op}}$ on a
closed subinterval of the given curve, then, for $|t|\le a_0$ and
$|t|b\le1$,
\begin{equation}\label{eq:curvesqlupper85}
 D_N(E,E(t))\le
 \min\left\{2,\;2b\sqrt N|t|+
             \tfrac12(L+16b^2)Nt^2\right\}.
\end{equation}
\end{lemma}
\begin{proof}
The factorization
$\mathcal C_E=T_A(I-AA^*)T_A^*/4$ in
Lemma~\ref{lem:covarianceidentities83} gives a horizontal solution:
for $K=\mathcal C_E^\dagger H$ take
$B=(I-AA^*)T_A^*K/4$.  Thus $T_AB=H$ and $A^*B=0$.
In fact $4\|B\|_{\mathrm{HS}}^2=
\langle H,\mathcal C_E^\dagger H\rangle$.
Formula \eqref{eq:horizontalcurve85} is normalized because
$(A+tB)^*(A+tB)=I+t^2S$.  Its derivative is
\[
 A'(t)=(B-tAS)(I+t^2S)^{-3/2},\qquad
 A'(t)^*A'(t)=S(I+t^2S)^{-2}.
\]
These identities prove horizontality and the first derivative bound.
Differentiation once more gives, in the displayed order,
\[
 A''(t)=-AS(I+t^2S)^{-3/2}
       -3t(B-tAS)S(I+t^2S)^{-5/2}.
\]
Since $\|A\|=1$, $\|S\|=b^2$, and $|t|b\le1$, its norm is at
most $b^2+3|t|(b+|t|b^2)b^2\le7b^2$.

Duplicate each output label into an inaccessible environment to form
a Stinespring isometry $W(t)$ from $A(t)$.  It has exactly the same
operator derivative bounds and $W(t)^*W'(t)=0$.
For one fixed purified adaptive tester, differentiating its final
vector gives a sum over slots.  Distinct terms are orthogonal: after
cancelling later common isometries, the later differentiated slot
contains $W^*W'$ or its adjoint.  Each diagonal term has squared
norm at most $b^2$.  The final vector derivative is at most
$\sqrt N b$, and the density derivative at most $2\sqrt N b$.
Integration followed by common partial trace proves
$D_N(E,\widetilde E(t))\le2b\sqrt N|t|$.
Padding stopped branches by discarded fixed-input calls includes
bounded public stopping; the original stopped output is a common
postprocessing of the padded experiment.

In the following application $\sum_jR_j=0$.  Thus the map is a
Hermiticity-preserving, trace-annihilating channel difference or
derivative, not a positive channel.  Lemma~\ref{lem:channelnorm86}
also proves its bound on arbitrary amplified trace-class inputs by
adjoint duality, without a state-optimization assumption.

For a Hermitian tuple $R$, the measurement derivative has diamond norm
at most $\sum_j\|R_j\|_{\mathrm{op}}$.
Indeed each reference output block is the partial trace against $R_j$;
its trace norm on a state is at most $\|R_j\|_{\mathrm{op}}$ by
duality.  The Hermiticity-preserving diamond norm may be optimized
over states with a reference.  The Stinespring second derivative
bound similarly gives
$\|\mathcal M_{\widetilde E(t)}''\|_\diamond
\le2\|W''(t)\|+2\|W'(t)\|^2\le16b^2$.
Taylor's integral formula, with common value and derivative at zero,
therefore bounds the one-call channel difference by
$(L+16b^2)t^2/2$.  Hybrid telescoping over $N$ calls proves
\eqref{eq:curvesqlupper85}.
\end{proof}

\begin{lemma}[A logical derivative for a general measurement curve]
\label{lem:curvelogical85}
Suppose $\Gamma=\Gamma_E(H)\notin\mathcal W_E$.
Use the support-complement code of Lemma~\ref{lem:supportcode84}
with $\Gamma$ in place of $G$, and let $\mathcal R$ be its fixed
CPTP recovery.  Put
\[
 \Phi_t=\mathcal R\circ(\mathcal M_{E(t)}\otimes\operatorname{Id})
             \circ\mathcal J,\qquad
 \Gamma_L=J^*(\Gamma\otimes I)J.
\]
Here $L$ is exactly the local curvature quantity
$\sup_{|u|\le a_0}\sum_j\|E_j''(u)\|_{\mathrm{op}}$ from
\eqref{eq:curvature86}; the remainder does not infer it from $E,H$.
Then $\Phi_0=\operatorname{Id}_2$, $\Phi'_0=-i[\Gamma_L,\cdot]$,
and, for $|t|\le a_0$,
\begin{equation}\label{eq:curvelogical85}
 \|\Phi_t-\operatorname{Ad}(e^{-it\Gamma_L})\|_\diamond
 \le\kappa t^2,\qquad
 \kappa=\tfrac12L+2\|\Gamma\|_{\mathrm{op}}^2.
\end{equation}
Its logical gap is
$\Delta=\|Z\|_{\mathrm{HS}}^2/\tr Z_+>0$,
where $Z=\Gamma-\Pi_{\mathcal W_E}\Gamma$.
\end{lemma}
\begin{proof}
To compute the derivative, not to factor the unknown curve itself,
put
\[
 A_j=\sqrt{E_j},\qquad
 B_j=(\sqrt{E_j})^+H_j(I-P_j/2).
\]
The inverse here is only on the fixed support; it is zero on its
complement.  Since $Q_jH_jQ_j=0$,
$A_j^*B_j+B_j^*A_j=H_j$ and
\[
 \sum_j A_j^*B_j=\tfrac12\sum_j[P_j,H_j]=-i\Gamma.
\]
In particular this sum is anti-Hermitian.  The normalized factors
$(A+tB)(I+t^2B^*B)^{-1/2}$ have derivative $B$ and induce the
same channel value and derivative as $E(t)$, irrespective of its
second-order rank changes.

Write $L_\mu=|j\rangle\langle a|A_j$ and
$\dot L_\mu=|j\rangle\langle a|B_j$ for $\mu=(j,a)$,
with identity on the reference understood.  If $R_\ell$ are Kraus
operators of the completed recovery, exact correction gives
$R_\ell L_\mu J=c_{\ell\mu}I_2$.
The left first-order term of the logical channel is therefore
\[
 \sum_{\ell,\mu}\overline{c_{\ell\mu}}
                  R_\ell\dot L_\mu J
 =J^*\sum_\mu L_\mu^*\dot L_\mu J
 =-i\Gamma_L,
\]
where completeness of the recovery is used in the first equality.
The adjoint term proves the commutator derivative.  The actual
logical second derivative is at most $L$ in diamond norm, and that
of the logical unitary is at most $4\|\Gamma\|_{\mathrm{op}}^2$.
Taylor's integral remainder proves \eqref{eq:curvelogical85}.
The support-code lemma gives the gap and the available systems.
\end{proof}

\begin{proof}[Proof of Theorem~\ref{thm:curveclassification85}]
If $H$ is in the covariance range, let $x=\sqrt N|t|$.
For $x\le1$, \eqref{eq:curvesqlupper85} is at most
$[2b+(L+16b^2)/2]x$; for $x\ge1$ use the bound two.
For the lower, choose $j,v$ with $\langle v,H_jv\rangle\ne0$.
The fixed input $v$ and the event ``label $j$'' have a first-order
probability gap, and Lemma~\ref{lem:bernoulli85} proves the result.

If $H$ is outside the range, apply Lemma~\ref{lem:curvelogical85}.
Starting from the equal logical superposition and telescoping $m$
corrected calls gives, for $|t|\le a_0$,
\begin{equation}\label{eq:curvefinitelower85}
 D_N(E,E(t))\ge
 \max_{1\le m\le N}\bigl(2|\sin(mt\Delta/2)|-m\kappa t^2\bigr)_+.
\end{equation}
Choose
\[
 t_0\le\min\{a_0,(2\Delta)^{-1},\Delta/(2\pi\kappa)\},\qquad
 m=\min\{N,\lfloor(|t|\Delta)^{-1}\rfloor\}.
\]
For $0<|t|\le t_0$ the floor is at least two, and
$x=m|t|\Delta$ satisfies
$\tfrac12\min\{1,N|t|\Delta\}\le x\le1$.
The sine term is at least $2x/\pi$ and the error at most $x/(2\pi)$.
This gives a positive constant times $\min\{1,N|t|\}$.
For the upper, the finite constant
$L_1=\sup_{|s|\le a_0}\sum_j\|E_j'(s)\|_{\mathrm{op}}$
and hybrid telescoping give $D_N\le\min\{2,NL_1|t|\}$.
The point $t=0$ is immediate.
\end{proof}

\begin{corollary}[Rank opening without a unitary orbit]\label{cor:rankopening85}
Let $P=|0\rangle\langle0|$ on $\C^2$,
$G=\left(\begin{smallmatrix}0&-i\\i&0\end{smallmatrix}\right)$,
$U_t=e^{itG}$, and for $|t|<1/2$ set
\[
 E_1(t)=(1-t^2)U_tPU_t^*+\tfrac12t^2I_2,
 \qquad E_2(t)=I_2-E_1(t).
\]
Both effects are full rank for $t\ne0$ and projective at zero.
Nevertheless $D_N(E(0),E(t))\asymp\min\{1,N|t|\}$ on a fixed
small interval, uniformly over $N\ge1$.
\end{corollary}
\begin{proof}
At zero the derivative is $i[G,(P,I-P)]\ne0$ and the support span
is diagonal, while $G$ is not.  Lemma~\ref{lem:curvegenerator85}
and Theorem~\ref{thm:curveclassification85} apply.  The eigenvalues
of the first effect are $t^2/2$ and $1-t^2/2$, proving the rank assertion.
\end{proof}
A vanishing first derivative is intentionally excluded from the
classification: it does not imply stationarity of a general curve.
For example the scalar curve $(t^2,1-t^2)$ has distance
$2[1-(1-t^2)^N]$ from $(0,1)$, of order $\min\{1,Nt^2\}$.
Likewise a one-sided rank opening can have a nonzero missing-support
first-order block and is not covered by the two-sided tangent lemma.

\section{Certified errors in the discrimination controls}
\label{sec:controlstability85}
The preceding lower bounds first concern ideal known-pair controls.
Their stability has a direct finite formulation.  An error bound on a
control channel is a bound in unhalved diamond norm, valid on every
input including its memory reference.  It is not a bound measured only
on one calibration state.

\begin{theorem}[A finite control-error ledger]\label{thm:controlstability85}
Use the code and recovery from Lemma~\ref{lem:curvelogical85}, with
gap $\Delta$ and remainder constant $\kappa$.
In this statement
$L=\sup_{|u|\le a_0}\sum_j\|E_j''(u)\|_{\mathrm{op}}$ and
$\kappa=L/2+2\|\Gamma_E(H)\|_{\mathrm{op}}^2$, as in
\eqref{eq:curvature86} and \eqref{eq:curvelogical85}.
Suppose an implemented initial logical state differs from the ideal
one in trace norm by at most $e_0$.  At cycle $r$, suppose its
implemented logical channel differs from the corresponding ideal
corrected channel by at most $\nu_r$ under each hypothesis.
Suppose the implemented final readout differs from the ideal
Helstrom readout by at most $e_f$ in diamond norm.
For every integer $m\le N$ and $|t|\le a_0$, the actual classical
output laws obey
\begin{equation}\label{eq:controlledger85}
 \|\widetilde p_{0,m}-\widetilde p_{t,m}\|_1
 \ge\left(2|\sin(mt\Delta/2)|-m\kappa t^2
              -2e_0-2e_f-2\sum_{r=1}^m\nu_r\right)_+.
\end{equation}
Here one may choose the Helstrom readout for the ideal unrotated and
logical-unitary states; it depends on the known pair and on $m$.
The cycle error may be bounded by the sum of the encoding and recovery
errors, because the unknown measurement itself is a channel.

In particular, put $z=\min\{1,N|t|\Delta\}$ and choose $m,t_0$
as in the proof of Theorem~\ref{thm:curveclassification85}.
For $0<|t|\le t_0$, the sufficient conditions
\begin{equation}\label{eq:controlprecision85}
 e_0+e_f\le z/(16\pi),\qquad
 \nu_r\le |t|\Delta/(8\pi)\quad(1\le r\le m)
\end{equation}
imply
$\|\widetilde p_{0,m}-\widetilde p_{t,m}\|_1\ge z/(2\pi)$.
All errors are deterministic uniform channel bounds; no independence
of control errors is required.
\end{theorem}
\begin{proof}
Choose the ideal two-outcome Helstrom readout for the equal logical
superposition and its image under $e^{-imt\Gamma_L}$.
Its classical $\ell^1$ separation is $2|\sin(mt\Delta/2)|$.
Replacing the logical unitary by the actual ideal corrected channel
loses at most $m\kappa t^2$, by channel telescoping and
Lemma~\ref{lem:curvelogical85}.  Under either hypothesis,
telescoping the implemented preparation, cycles and readout changes
its output law in $\ell^1$ by at most
$e_0+e_f+\sum_r\nu_r$.  The triangle inequality proves
\eqref{eq:controlledger85}.  Replacing encoding and recovery around
one unchanged unknown channel costs no more than the sum of their
diamond errors, which proves the stated decomposition of $\nu_r$.

For the indicated choice set $x=m|t|\Delta$, so $z/2\le x\le1$.
The ideal sine term minus curvature is at least $3x/(2\pi)$.
Conditions \eqref{eq:controlprecision85} give
$e_0+e_f\le x/(8\pi)$ and $\sum_r\nu_r\le x/(8\pi)$.
The total loss is therefore at most $x/(2\pi)$, leaving
$x/\pi\ge z/(2\pi)$.
\end{proof}

For a unitary orbit, Lemma~\ref{lem:correctedchannel84} gives the
same result with $\kappa=4\|G\|_{\mathrm{op}}^2$ and the original
support-code gap.  An error allowance fixed independently of the
angle need not meet \eqref{eq:controlprecision85} as $t\to0$.
The theorem does not certify a particular piece of hardware, synthesize
a gate circuit, or allow target-dependent advice in common learning.
It specifies exactly which certified approximation errors suffice
to preserve the finite-use lower bound.

\begin{proposition}[Exact first-order decisions]\label{prop:curveexact85}
For Gaussian-rational $E$ and a Gaussian-rational Hermitian zero-sum
$H$, two-sided first-order realizability and the branch of
Theorem~\ref{thm:curveclassification85} for $H\ne0$ are decidable
in polynomial represented-input bit complexity.  The decision gives
$\Gamma_E(H)$, its projection on $\mathcal W_E$, and its orthogonal
residual.  A zero tangent is reported as a higher-order question,
not as stationarity of an unspecified curve.
\end{proposition}
\begin{proof}
The exact support projectors in Proposition~\ref{prop:supportexact84}
are rational.  Test $Q_jH_jQ_j=0$ for every $j$ and compute
\eqref{eq:curvegenerator85}.  Necessity and sufficiency of the test
follow from Lemma~\ref{lem:curvegenerator85} and the normalized-factor
construction in Lemma~\ref{lem:curvelogical85}: for any such tangent
it gives a real-analytic measurement curve with that derivative.
Real coordinate rank selection and the Hilbert--Schmidt Gram system
then project $\Gamma_E(H)$ exactly.  Rational matrix arithmetic and
fraction-free elimination have the bit bounds already proved in
Proposition~\ref{prop:supportexact84}.  Only these finite operations
are performed.  A numerical curvature bound for an unspecified
curve, an optimal tester, and a compiled recovery are not outputs
of this decision.
\end{proof}

\section{One-sided tangent neighborhoods}
\label{sec:tangentneighborhood86}
The preceding curve results admit a finite-pair formulation.  The
base measurement and its direction are fixed, but the second-order
perturbation may vary arbitrarily within a prescribed bound.  This
also permits first-order opening of a missing support.  Write
$\|R\|_{\Sigma}=\sum_j\|R_j\|_{\mathrm{op}}$ for a Hermitian tuple.
All trace and diamond norms in this section are unhalved.

\begin{lemma}[A reference-uniform channel bound]\label{lem:channelnorm86}
For a Hermitian tuple $R$, let
$\mathcal D_R(X)=\sum_j\tr(R_jX)|j\rangle\langle j|$.
Then
\begin{equation}\label{eq:channelnorm86}
 \|\mathcal D_R\|_\diamond\le\|R\|_{\Sigma}.
\end{equation}
The map is Hermiticity preserving, and is trace annihilating when
$\sum_jR_j=0$.  In particular this applies to measurement-channel
differences, derivatives, and Taylor remainders.
\end{lemma}
\begin{proof}
At every reference dimension the output is block diagonal.  In trace
norm duality it therefore suffices to use a block-diagonal test
$Y=\bigoplus_jY_j$, with $\|Y_j\|_{\mathrm{op}}\le1$.  Its adjoint
image is $\sum_jR_j\otimes Y_j$, of norm at most
$\sum_j\|R_j\|_{\mathrm{op}}$.  Duality bounds the amplified map on
\emph{every} trace-class input $X$, not only on states.  Taking the
supremum over reference dimensions proves \eqref{eq:channelnorm86}.
The remaining assertions follow by taking adjoints and traces.
\end{proof}

\begin{lemma}[The one-sided cone and its lineality]\label{lem:onesidedcone86}
Fix a measurement $E$, put $P_j=\operatorname{supp}E_j$ and $Q_j=I-P_j$,
and let $H$ be a Hermitian zero-sum tuple.  It is the right derivative
at zero of a one-sided real-analytic measurement curve if and only if
\begin{equation}\label{eq:onesidedcone86}
 J_j:=Q_jH_jQ_j\succeq0\quad\hbox{for every }j.
\end{equation}
The largest vector subspace of this cone is given by $J_j=0$ for all
$j$.  For every Hermitian zero-sum tuple, not necessarily in the cone,
\begin{equation}\label{eq:fullrange86}
 H\in\operatorname{Ran}\mathcal C_E
 \quad\Longleftrightarrow\quad
 J_j=0\ (1\le j\le k),\qquad
 \Gamma_E(H)\in\mathcal W_E.
\end{equation}
Here $\Gamma_E(H)=\frac i2\sum_j[P_j,H_j]$ and
$\mathcal W_E=\sum_jP_j\mathcal H_dP_j$ are real Hermitian objects.
\end{lemma}
\begin{proof}
Right differentiation of $\langle v,E_j(t)v\rangle\ge0$ for
$v\in\operatorname{Ran}Q_j$ proves necessity.  For sufficiency, for
each nonzero support let $\lambda_j>0$ be the smallest eigenvalue
of $E_j$ on that support.  Choose
\[
 c\ge1+\max_{j:P_j\ne0}2\|H_j\|_{\mathrm{op}}^2/\lambda_j.
\]
For all sufficiently small $s\ge0$, the $P_j$ block of
$E_j+sH_j+cs^2I$ is at least $\lambda_jP_j/2$.
Its Schur complement on $Q_j$ is bounded below by
\[
 sJ_j+s^2\left(c-2\|H_j\|_{\mathrm{op}}^2/\lambda_j\right)Q_j
 \succeq0.
\]
If $P_j=0$, condition \eqref{eq:onesidedcone86} is $H_j\succeq0$
and positivity is immediate.  Thus the rational matrix curve
\begin{equation}\label{eq:conecurve86}
 E_j^{\mathrm c}(s)=\frac{E_j+sH_j+cs^2I}{1+kcs^2}
\end{equation}
is positive, sums exactly to $I$, and has right derivative $H_j$.
This proves sufficiency, including singular $J_j$ with nonzero cross
blocks.  The cone and its negative intersect precisely at $J_j=0$.

Pair $H$ with the complete kernel in
Theorem~\ref{thm:supportkernel84}.  Taking only its independent
missing-support parameters $Z_j$ shows that orthogonality to the
kernel forces $Q_jH_jQ_j=0$; the tuple-mean term vanishes because
$\sum_jH_j=0$.  Once those blocks vanish, the remaining pairing is
$-2\tr(A\Gamma_E(H))$ for every $A\in\mathcal W_E^\perp$.
Self-adjointness of $\mathcal C_E$ proves \eqref{eq:fullrange86}.
\end{proof}

\begin{lemma}[A finite normalized-factor remainder]\label{lem:factorremainder86}
Let $A_j^*A_j=E_j$, $A^*A=I$, and
$A_j^*B_j+B_j^*A_j=H_j$, with $\sum_jH_j=0$.
Set $b=\|B\|_{\mathrm{op}}$, $h=\max_j\|H_j\|_{\mathrm{op}}$ and
\[
 R_s=(I+s^2B^*B)^{-1/2},\qquad
 G_j(s)=R_s(A_j+sB_j)^*(A_j+sB_j)R_s.
\]
This defines an analytic measurement near zero, and, for $0\le s\le1$,
\begin{equation}\label{eq:factorremainder86}
 \|G(s)-E-sH\|_{\Sigma}\le K_b s^2,
 \qquad K_b=kb^2(2+h).
\end{equation}
When $sb\le1$, its stacked factors differ from $A$ in operator norm
by at most $3bs/2$.  If $A^*B=0$, then also
\begin{equation}\label{eq:horizontalpair86}
 D_N(E,G(s))\le\min\{2,2b\sqrt N s\}.
\end{equation}
\end{lemma}
\begin{proof}
The zero tuple sum gives $A^*B+B^*A=0$, so the stacked normalization
is $I+s^2B^*B$.  This matrix is positive definite for all real $s$;
its inverse square root is analytic near zero.  Functional calculus
gives $\|R_s\|\le1$ and $\|R_s-I\|\le s^2b^2/2$.
Expanding the effect as
$R_s(E_j+sH_j+s^2B_j^*B_j)R_s$ bounds its remainder by
$s^2b^2(1+sh)+s^2\|B_j\|^2$.  Sum over $j$, use
$\sum_j\|B_j\|^2\le kb^2$, and obtain \eqref{eq:factorremainder86}.
The factor difference is at most $sb+s^2b^2/2$.
For horizontal $B$, the ordered identities in
Lemma~\ref{lem:horizontalcurve85} give $A(s)^*A'(s)=0$ and
$\|A'(s)\|\le b$.  Copying the classical label into a discarded
environment is an isometric embedding of the stacked factors, so
it preserves these derivative norms.  The orthogonal-slot argument
of that lemma gives \eqref{eq:horizontalpair86}, including bounded
public stopping.  The environment is not supplied to the tester.
\end{proof}

\begin{definition}[A quadratic tangent neighborhood]\label{def:tangenttube86}
Fix $E$, a Hermitian zero-sum $H\ne0$ satisfying
\eqref{eq:onesidedcone86}, and a number $\Lambda\ge0$.
For $s>0$, the neighborhood consists of all legal measurements $F$
with
\begin{equation}\label{eq:tangenttube86}
 \|F-E-sH\|_{\Sigma}\le\Lambda s^2.
\end{equation}
The number $\Lambda$ is a bound on a \emph{finite remainder}, not an
inferred curvature bound.  No curve, or differentiability of $F$ as
a function of $s$, is part of this definition.
\end{definition}

\begin{theorem}[Finite-pair classification in a fixed tangent neighborhood]
\label{thm:tubedichotomy86}
For each fixed $E,H,\Lambda$ in Definition~\ref{def:tangenttube86}
there are $c,C,s_0>0$ such that, for every integer $N\ge1$, every
$0<s\le s_0$, and every legal $F$ satisfying
\eqref{eq:tangenttube86},
\begin{equation}\label{eq:tubedichotomy86}
 c\min\{1,N^\alpha s\}\le D_N(E,F)
       \le C\min\{1,N^\alpha s\},\qquad
 \alpha=\begin{cases}
  1/2,&H\in\operatorname{Ran}\mathcal C_E,\\
  1,&H\notin\operatorname{Ran}\mathcal C_E.
 \end{cases}
\end{equation}
The constants are uniform over the stated remainder neighborhood,
but not over $E,H$ or $\Lambda$.  If some $J_j\ne0$, the linear
lower is attained by repeated product inputs and a classical event.
If all $J_j=0$ but $\Gamma_E(H)\notin\mathcal W_E$, it is attained
by the known-direction correction code on the actual label and a
retained reference, with ideal controls.  Neither construction is
a common unknown-measurement learner.
\end{theorem}

There are two different reasons for a linear scale.  First-order
support opening creates an event impossible at the base measurement.
Within the lineality space of the cone, coherent correction produces
the established Kraus-span metrological mechanism \cite{ZJ84,KL84}.
The result above makes this distinction in support coordinates and
controls finite remainders uniformly at a fixed base and direction.
It is not a uniform arbitrary-pair equivalence for the midpoint
covariance certificate.

\begin{proof}
Put $h=\max_j\|H_j\|_{\mathrm{op}}>0$ and
$h_1=\|H\|_{\Sigma}$.  Choose a unit eigenvector of one $H_j$
with expectation of magnitude $h$.  On that fixed input the event
``label $j$'' has probability gap at least $hs-\Lambda s^2$.
For $s$ sufficiently small this is at least $hs/2$.
Lemma~\ref{lem:bernoulli85} therefore gives a uniform product lower
of order $\min\{1,\sqrt N s\}$, for every $F$ in the neighborhood.
All pairs also satisfy the elementary upper
\begin{equation}\label{eq:tubehybrid86}
 D_N(E,F)\le\min\{2,N(h_1s+\Lambda s^2)\},
\end{equation}
by Lemma~\ref{lem:channelnorm86} and channel telescoping.

Suppose first that $H\in\operatorname{Ran}\mathcal C_E$.
The horizontal solution of Lemma~\ref{lem:horizontalcurve85}
provides factors $B$ with $A^*B=0$ and $T_AB=H$.
Lemma~\ref{lem:factorremainder86} and \eqref{eq:tangenttube86} give
\begin{equation}\label{eq:tubefiniteupper86}
 D_N(E,F)\le\min\{2,2b\sqrt N s+(\Lambda+K_b)Ns^2\}.
\end{equation}
For $x=\sqrt N s\le1$, the right side is at most
$(2b+\Lambda+K_b)x$; for $x\ge1$ it is at most two.
This proves the upper in the first branch without discarding its
finite $Ns^2$ remainder before comparison.

Next suppose some $J_j\ne0$.  It is positive semidefinite; choose a
unit vector $v\in\operatorname{Ran}Q_j$ with
$\lambda=\langle v,J_jv\rangle>0$.
The base probability of label $j$ is zero, while that under $F$ is
at least $\lambda s-\Lambda s^2\ge\lambda s/2$ after decreasing
$s_0$.  Repeating $v$ gives the exact impossible-event bound
\begin{equation}\label{eq:leakagelower86}
 D_N(E,F)\ge2\left[1-(1-\lambda s/2)^N\right]
 \ge2(1-e^{-1})\min\{1,N\lambda s/2\}.
\end{equation}
Here $s_0$ is chosen also so that $\lambda s/2\le1$.
Together with \eqref{eq:tubehybrid86} this proves the linear law.
The vector and the decision event are the same for every allowed $F$.

It remains to consider $J_j=0$ for all $j$ and
$\Gamma:=\Gamma_E(H)\notin\mathcal W_E$.
Use the completed recovery and logical code of
Lemma~\ref{lem:supportcode84}, applied to $\Gamma$.
The normalized first-order factors in Lemma~\ref{lem:curvelogical85}
show algebraically that the corrected derivative $\mathcal D_H$
is $-i[\Gamma_L,\cdot]$.  This calculation needs only $E,H$, not
a smooth factorization of $F$.
For $\Phi_F=\mathcal R\circ(\mathcal M_F\otimes\operatorname{Id})
\circ\mathcal J$, \eqref{eq:tangenttube86} gives
\[
 \|\Phi_F-\operatorname{Id}_2+is[\Gamma_L,\cdot]\|_\diamond
 \le\Lambda s^2.
\]
The integral Taylor remainder for the logical unitary is at most
$2\|\Gamma\|_{\mathrm{op}}^2s^2$.  Consequently
\begin{equation}\label{eq:tubelogical86}
 \|\Phi_F-\operatorname{Ad}(e^{-is\Gamma_L})\|_\diamond
 \le\kappa_{\mathrm{tube}}s^2,
 \qquad \kappa_{\mathrm{tube}}=\Lambda+2\|\Gamma\|_{\mathrm{op}}^2.
\end{equation}
The gap is $\Delta=\|Z\|_{\mathrm{HS}}^2/\tr Z_+>0$, where
$Z=\Gamma-\Pi_{\mathcal W_E}\Gamma$.
For every $m\le N$, equal logical amplitudes and channel telescoping
give a lower bound
$\bigl(2|\sin(ms\Delta/2)|-m\kappa_{\mathrm{tube}}s^2\bigr)_+$.
Choose
\[
 s_0\le\min\{1,(2\Delta)^{-1},
              \Delta/(2\pi\kappa_{\mathrm{tube}})\},\qquad
 m=\min\{N,\lfloor(s\Delta)^{-1}\rfloor\}.
\]
The floor is at least two, and
$\frac12\min\{1,Ns\Delta\}\le x=ms\Delta\le1$.
The sine term is at least $2x/\pi$ and the loss at most $x/(2\pi)$.
In particular the separation is at least
$\min\{1,Ns\Delta\}/(2\pi)$.
The readout may be fixed to the ideal logical pair, so this lower is
uniform over all the permitted $F$ as well.  Its code uses an external
reference of dimension at most $2d$ per call.  Finally
\eqref{eq:tubehybrid86} completes the proof.
\end{proof}

\begin{corollary}[All nonzero one-sided first jets]\label{cor:onesidedcurves86}
Let $E(t)$, $0\le t<a$, be a fixed $C^2$ measurement curve, including its
right derivatives at zero, with $H=E'(0)\ne0$.  Then
\eqref{eq:tubedichotomy86} holds with $s=t$ and curve-dependent
constants on a fixed right interval.  The branch is decided by
\eqref{eq:fullrange86}.  Thus first-order support openings, as well
as tangential second-order rank changes, are included.
\end{corollary}
\begin{proof}
One-sided positivity gives \eqref{eq:onesidedcone86}.
On $[0,a_0]$ set
$L_+=\sup_{0\le u\le a_0}\sum_j\|E_j''(u)\|_{\mathrm{op}}$.
Taylor's formula gives \eqref{eq:tangenttube86} with
$\Lambda=L_+/2$.  Apply Theorem~\ref{thm:tubedichotomy86}.
This does not extend the two-sided tangent lemma by dropping one of
its hypotheses; it uses the larger cone of Lemma~\ref{lem:onesidedcone86}.
\end{proof}

\begin{corollary}[Certified controls on the whole neighborhood]
\label{cor:tubecontrols86}
In the tangential linear branch, Theorem~\ref{thm:controlstability85}
holds uniformly over \eqref{eq:tangenttube86} with $t=s$ and
$\kappa=\kappa_{\mathrm{tube}}$ from \eqref{eq:tubelogical86}.
In particular the same initial, cycle and final error allowances
with $z=\min\{1,Ns\Delta\}$ preserve separation at least $z/(2\pi)$.
\end{corollary}
\begin{proof}
Equation~\eqref{eq:tubelogical86} supplies the only curvature input
used in that theorem's telescoping proof.  Pay it once, and pay the
certified implementation errors under both hypotheses exactly as
in \eqref{eq:controlledger85}.  The ideal readout depends on
$E,H,s,m$, not on the unknown remainder within the neighborhood.
The precision allowances remain angle dependent, not fixed-noise
robustness or a control-synthesis guarantee.
\end{proof}

\section{Independent probes and higher-order approaches}
\label{sec:independent86}
A linear finite-use scale need not signal coherent accumulation.
The support-opening case can already be distinguished by independent
classical observations.  We make the acquisition restriction precise.
Let $D_N^{\mathrm{prod}}(E,F)$ be the supremum of the final trace
separation when the input is a product of at most $N$ probe--reference
states, one pair per call.  The pairs may differ and may depend on the
known hypotheses.  All completed outputs and references may undergo
an arbitrary joint final readout.  There is no quantum correlation
between different input pairs, and no intervening feedback into later
inputs.  Public randomization is allowed.  This is narrower than
nonadaptive acquisition with an entangled many-call input.

\begin{theorem}[Three local acquisition regimes]\label{thm:producttrichotomy86}
Fix $E,H,\Lambda$ as in Definition~\ref{def:tangenttube86}.  Uniformly
in $N$, all sufficiently small $s>0$, and every $F$ in
\eqref{eq:tangenttube86}, the two distances have the following orders:
\begin{center}
\begin{tabular}{@{}lll@{}}
\toprule
Direction & $D_N^{\mathrm{prod}}(E,F)$ & $D_N(E,F)$\\
\midrule
$J_j=0$ for all $j$, $\Gamma_E(H)\in\mathcal W_E$
 & $\min\{1,\sqrt N s\}$ & $\min\{1,\sqrt N s\}$\\
$J_j=0$ for all $j$, $\Gamma_E(H)\notin\mathcal W_E$
 & $\min\{1,\sqrt N s\}$ & $\min\{1,Ns\}$\\
Some $J_j\ne0$
 & $\min\{1,Ns\}$ & $\min\{1,Ns\}$\\
\bottomrule
\end{tabular}
\end{center}
Every comparison constant and the common small interval depend only
on the fixed $E,H,\Lambda$.  In particular, collective processing of
independently acquired outputs does not recover the linear scale of
the second row.  The assertion does not bound arbitrary entangled
parallel inputs or all classical-feedback acquisition models.
\end{theorem}
\begin{proof}
Theorem~\ref{thm:tubedichotomy86} gives the adaptive column.  Its
repeated-input lower gives the product lower in the first two rows;
\eqref{eq:leakagelower86} and \eqref{eq:tubehybrid86} give both
product bounds in the last row.  It remains to prove a
reference-uniform product upper when all $J_j=0$.

Use the canonical first-order factors
$A_j=\sqrt{E_j}$ and $B_j=(\sqrt{E_j})^+H_j(I-P_j/2)$ from
Lemma~\ref{lem:curvelogical85}.  They need not be horizontal.
Put $b=\|B\|$, let $G(s)$ be their normalized measurement in
Lemma~\ref{lem:factorremainder86}, and set
\[
 a=\tfrac32 b+\sqrt{\Lambda+K_b}.
\]
For every input state with any reference, denote the corresponding
output states under $E,G(s),F$ by $\rho_0,\rho_g,\rho_f$.
Their Bures distance
$\beta(\rho,\sigma)=\sqrt{2-2f(\rho,\sigma)}$, where
$f(\rho,\sigma)=\|\sqrt\rho\sqrt\sigma\|_1$ is root fidelity,
satisfies
\[
 \beta(\rho_0,\rho_g)\le\tfrac32bs,
 \qquad
 \beta(\rho_g,\rho_f)^2\le\|\rho_g-\rho_f\|_1
                        \le(\Lambda+K_b)s^2.
\]
The first inequality uses purifications from the two stacked factors
and their operator-norm difference; it holds uniformly over the input.
The second uses \eqref{eq:channelnorm86},
\eqref{eq:factorremainder86}, and the fidelity--trace inequalities.
The triangle inequality for $\beta$ follows by choosing common
purifications at the middle state and applying the Hilbert-space
triangle inequality.  Thus $1-f(\rho_0,\rho_f)\le a^2s^2/2$.
These standard fidelity properties and their unhalved trace convention
are as in \cite[Proposition~3.16 and Theorems~3.22, 3.33]{Watrous65}.

For different independent probe--reference pairs the output states
are still products.  Root fidelity is multiplicative, and
$1-\prod_r f_r\le\sum_r(1-f_r)$.  Therefore
\[
 \left\|\bigotimes_r\rho_{0,r}-\bigotimes_r\rho_{f,r}\right\|_1
 \le2\sqrt{1-\prod_r f_r^2}
 \le2a\sqrt N s.
\]
A joint final readout contracts trace norm.  Convexity gives the same
bound for public randomization, including a retained record of that
randomization.  Capping the result at two proves the required upper.
\end{proof}

\begin{corollary}[Higher-order approaches with a quadratic remainder]
\label{cor:powerjets86}
Let $q\ge1$ be an integer, and let a legal measurement family $F(t)$,
$0<t<a$, satisfy for fixed $E,H\ne0,\Lambda$
\begin{equation}\label{eq:powerjet86}
 \|F(t)-E-t^qH\|_{\Sigma}\le\Lambda t^{2q}.
\end{equation}
Then $H$ satisfies \eqref{eq:onesidedcone86}, and the conclusions of
Theorems~\ref{thm:tubedichotomy86} and
\ref{thm:producttrichotomy86} hold with $s=t^q$.
This includes zero first derivatives when $q>1$.
\end{corollary}
\begin{proof}
Compress \eqref{eq:powerjet86} to each missing support, divide by
$t^q$, and let $t\downarrow0$; positivity gives $J_j\succeq0$.
Substitute $s=t^q$ in the two finite-pair theorems.  No inverse
reparametrization is differentiated.  In particular, the conclusion
holds for $F(t)=E_*(t^q)$ whenever $E_*$ is a fixed one-sided $C^2$
curve with nonzero right derivative.  It also holds for a $C^{2q}$
curve whose first nonzero derivative has order $q$ and whose
intermediate derivatives of orders $q+1,\ldots,2q-1$ vanish, by Taylor's remainder after the leading nonzero coefficient.
\end{proof}

Condition \eqref{eq:powerjet86} is substantive.  A first nonzero
coefficient alone does not imply it.  The following exact family
displays what an intervening support-opening term can contribute.

\begin{proposition}[Mixed higher-order rates]\label{prop:mixedrates86}
Let $1\le a<b<2a$ be integers.  For scalar input dimension and three
outcomes put
\[
 E=(\tfrac12,\tfrac12,0),\qquad
 F(t)=\bigl((1-t^b)(\tfrac12+t^a),
            (1-t^b)(\tfrac12-t^a),t^b\bigr),
 \quad 0<t\le\tfrac14.
\]
There are absolute positive constants $c,C$ such that for all $N\ge1$
\begin{equation}\label{eq:mixedrates86}
 c\min\{1,\sqrt N t^a+Nt^b\}
 \le D_N(E,F(t))=D_N^{\mathrm{prod}}(E,F(t))
 \le C\min\{1,\sqrt N t^a+Nt^b\}.
\end{equation}
The leading coefficient $(1,-1,0)$ is in the covariance range.
Its order-$a$ contribution and the order-$b$ opening must both be
retained to obtain the finite-use law.
\end{proposition}
\begin{proof}
A scalar-input device only produces independent classical draws.
Every tester output, including one with bounded stopping, is a common
postprocessing of the full $N$-draw record.  Thus its optimal distance
is exactly the $\ell^1$ distance of the two product laws.
The event that label three occurs gives the lower
$2[1-(1-t^b)^N]\ge2(1-e^{-1})\min\{1,Nt^b\}$.
Send labels one and two to their corresponding bits and label three
to a fair bit.  This gives Bernoulli parameters $1/2$ and
$1/2+(1-t^b)t^a$; their gap is at least $t^a/2$.
Lemma~\ref{lem:bernoulli85} gives a second lower of order
$\min\{1,\sqrt N t^a\}$.  The larger of these lower bounds controls
one half of the truncated sum.

Under $F(t)$ the probability of no third label is $(1-t^b)^N$;
conditioned on that event the record is Bernoulli with parameter
$1/2+t^a$.  Splitting off the complementary event gives the upper
\[
 2[1-(1-t^b)^N]
 +\|\operatorname{Ber}(1/2)^{\otimes N}
       -\operatorname{Ber}(1/2+t^a)^{\otimes N}\|_1.
\]
For $u=t^a$, the one-copy root fidelity obeys
$f^2=(1+\sqrt{1-4u^2})/2$, so $1-f^2\le2u^2$.
Multiplicativity and the trace upper give a bound $2\sqrt{2N}\,u$
for the Bernoulli term.  Use $1-(1-t^b)^N\le Nt^b$ and cap at two.
For the range assertion, the scalar support generator is zero and
the missing-support component of $(1,-1,0)$ vanishes; apply
\eqref{eq:fullrange86}.  Since $b<2a$, the opening remainder is not
$O(t^{2a})$, so this example does not satisfy
\eqref{eq:powerjet86} with that leading coefficient.
\end{proof}

For example, $a=2,b=3$ gives
$\min\{1,\sqrt N t^2+Nt^3\}$ despite a vanishing first derivative.
The scalar family $(t^q,1-t^q)$, in contrast, has the exact distance
$2[1-(1-t^q)^N]$ and satisfies \eqref{eq:powerjet86} in the opening
branch.  A zero tangent supplied without higher-order data is still
an undetermined higher-order problem, not a statement of stationarity.

\begin{proposition}[Exact cone and finite-neighborhood certificates]
\label{prop:coneexact86}
For Gaussian-rational $E,H$, one-sided first-order realizability,
lineality, covariance-range membership and the three mechanisms of
Theorem~\ref{thm:producttrichotomy86} are decidable in polynomial
represented-input bit complexity.  For rational $s>0$, $\Lambda\ge0$
and $\lambda_j\ge0$ with $\sum_j\lambda_j\le\Lambda$, legality of
$F$ and the sufficient finite-neighborhood inequalities
\begin{equation}\label{eq:pairpredicate86}
 -\lambda_js^2I\preceq F_j-E_j-sH_j\preceq\lambda_js^2I
 \quad(1\le j\le k)
\end{equation}
are also exactly decidable with that complexity.
A nonzero opening supplies a rationally represented witness and its
exact output probability at a supplied pair.  No local interval,
curvature bound, recovery circuit or exact adaptive distance is
inferred from a first jet alone.
\end{proposition}
\begin{proof}
The rational support projectors of
Proposition~\ref{prop:supportexact84} give every $J_j$.
Exact positive-semidefinite tests establish \eqref{eq:onesidedcone86};
zero tests establish lineality.  When the $J_j$ vanish, the same
real-rank and Hilbert--Schmidt Gram projection decides whether
$\Gamma_E(H)$ belongs to $\mathcal W_E$.
Lemma~\ref{lem:onesidedcone86} gives necessity and sufficiency,
not only a formal positivity prefilter.

If $J_j\succeq0$ is nonzero, some diagonal element $(J_j)_{rr}$ is
positive.  Put $v=Q_je_r$.  Then $v\ne0$,
$v^*H_jv=(J_j)_{rr}>0$, and $\|v\|^2=(Q_j)_{rr}>0$.
Both the rate $v^*H_jv/\|v\|^2$ and, for a supplied $F$,
the probability $v^*F_jv/\|v\|^2$ are rational.
The unit vector itself need not have rational coordinates: its
rational vector and squared normalization specify the state exactly.
Finally \eqref{eq:pairpredicate86} is a finite list of rational PSD
tests.  Rational elimination, support projection and these tests
have the polynomial bit bounds of the retained exact kernels.
Verifying one supplied pair does not verify a neighborhood for an
unspecified curve, or that a supplied $s$ lies in the theorem's
small interval.  The exact impossible-event probability, when
provided, has its finite $N$-draw interpretation without that caveat.
\end{proof}

\section{Regularization transported by the output channel}\label{sec:ridge84}
The fixed Euclidean ridge in Theorem~\ref{thm:closedcovariance83}
is tied to the displayed outcome coordinates.  Its contraction under
arbitrary stochastic relabeling does not follow from the
unregularized data-processing theorem.  A ridge obtained from a
public probability vector has a different, covariant property.

For a probability vector $w=(w_j)$, put $W=(w_jI_d)_j$ and
$\mathcal B_w=\mathcal C_W$.  On the full real Hermitian tuple
space this operator is
\begin{equation}\label{eq:weightedridge84}
 (\mathcal B_wK)_j=w_j\left(K_j-\sum_\ell w_\ell K_\ell\right).
\end{equation}
For a zero-sum direction $H$ and $\tau>0$ define the extended energy
\begin{equation}\label{eq:transportenergy84}
 \mathcal E_{E,w,\tau}(H)=\sup_{K\in\mathcal H_d^k}
 \left\{2\langle H,K\rangle-
       \langle K,(\mathcal C_E+\tau\mathcal B_w)K\rangle\right\}.
\end{equation}
It equals the pseudoinverse quadratic form on the range and is
$+\infty$ off the range.  Constant tuples have zero energy and
zero pairing with $H$.  If every $w_j>0$, the supremum is finite
and may be evaluated by restricting to $\mathcal T$.

\begin{theorem}[Processing-covariant regularized energy]\label{thm:transportedridge84}
For every column-stochastic output map $T$, every $\tau>0$, and
every zero-sum Hermitian tuple $H$,
\begin{equation}\label{eq:ridgeprocessing84}
 \mathcal E_{TE,Tw,\tau}(TH)\le\mathcal E_{E,w,\tau}(H).
\end{equation}
Zero entries of $Tw$ are allowed, with the extended-energy convention
in \eqref{eq:transportenergy84}.  If a stochastic map $S$ satisfies
$ST=\operatorname{Id}$, equality holds in
\eqref{eq:ridgeprocessing84}.  In particular, reversible outcome
splitting leaves this energy unchanged when the public weights
are split by the same rule.
\end{theorem}
\begin{proof}
The matrix Jensen identity in
Proposition~\ref{prop:covarianceprocessing83} is an inequality on
the full tuple space, not only on its zero-sum subspace.  It gives
\[
 \langle L,\mathcal C_{TE}L\rangle
 \ge\langle T^*L,\mathcal C_E T^*L\rangle.
\]
Apply the same inequality to the scalar measurement $W$ to get
$\langle L,\mathcal B_{Tw}L\rangle
\ge\langle T^*L,\mathcal B_wT^*L\rangle$.
Also $\langle TH,L\rangle=\langle H,T^*L\rangle$.
Thus each expression in the supremum defining the left side of
\eqref{eq:ridgeprocessing84} is bounded by a permissible expression
on the right.  Taking suprema proves the claim, including infinite
values.  No assumption that $T^*L$ has zero sum is needed.
If $ST=\operatorname{Id}$, application of the same inequality to
$S$ gives the reverse bound.
\end{proof}

\begin{proposition}[Relation to the operational upper certificate]\label{prop:ridgeupper84}
If $w_j>0$ and $w_{\max}=\max_jw_j$, then for every pair of legal
measurements, with midpoint $M$ and difference $H$,
\begin{equation}\label{eq:ridgeupper84}
 D_N(E,F)\le\min\left\{2,
 2\sqrt{k+1}\sqrt{N\mathcal E_{M,w,1/(Nw_{\max})}(H)}\right\}.
\end{equation}
For $w_j=1/k$, one has $\mathcal B_w|_{\mathcal T}
=\operatorname{Id}/k$, and
$\mathcal E_{M,w,k/N}(H)=\mathcal Q_N(E,F)^2/N$ exactly.
\end{proposition}
\begin{proof}
The variance expression gives
$0\preceq\mathcal B_w\preceq w_{\max}\operatorname{Id}$.
On $\mathcal T$, all coordinates of $w$ positive make
$\mathcal B_w$ positive definite.  Inverse order therefore gives
\[
 \langle H,(\mathcal C_M+N^{-1}\operatorname{Id})^{-1}H\rangle
 \le\mathcal E_{M,w,1/(Nw_{\max})}(H).
\]
Use Theorem~\ref{thm:closedcovariance83}.  The uniform formula
follows by substitution into \eqref{eq:weightedridge84}.
\end{proof}
Theorem~\ref{thm:transportedridge84} compares the same $\tau$ and
\emph{transported} weights.  Replacing $Tw$ by a new uniform vector,
or changing $\tau$ to the new alphabet size divided by $N$, is a
different operation.  Consequently this result does not retroactively
assert arbitrary-processing monotonicity of the original fixed ridge.
Exact evaluation with positive rational weights uses the tangent Gram
system of Proposition~\ref{prop:covarianceexact83}, with the ridge
stiffness $\tau\mathcal B_w$ in place of $G/N$.

\section{Finite-outcome geometry and exact descriptions}
\label{sec:finiteoutcome82}

The interior geometry does not require two outcomes.  We now allow any
fixed finite alphabet, while retaining the input-consuming, ordered,
classical-output interface.  The effects below need not commute.
The binary results on the complete effect body remain separate: no
positive lower spectral bound is imposed in those results.

For integers $d\ge1$ and $k\ge2$, put
\begin{equation}\label{eq:balancedbody82}
 \mathfrak P_{d,k}=\left\{\boldsymbol E=(E_1,\ldots,E_k):
 E_j=E_j^*,\quad \sum_{j=1}^k E_j=I_d,\quad
 \frac{I_d}{2k}\preceq E_j\preceq\frac{3I_d}{2k}\right\}.
\end{equation}
An element specifies the channel
$\mathcal M_{\boldsymbol E}(\rho)
=\sum_j\operatorname{tr}(E_j\rho)|j\rangle\langle j|$.
It returns no residual quantum system.  Write $D_N$ for its
$N$-use distance, with precisely the reference, feedback, bounded
public stopping, and unhalved trace-norm conventions defining $d_N$.
The superscript ${\rm na}$ denotes nonadaptive experiments, not a
restriction to unentangled inputs.  Set
\[
 r(\boldsymbol E,\boldsymbol F)=\max_j\|E_j-F_j\|_{\rm op},
 \qquad p=(k-1)d^2.
\]
The affine space $\sum_j E_j=I_d$ has real dimension $p$.
In its translated vector space $\sum_j H_j=0$, the set
$\mathfrak P_{d,k}$ is exactly the closed $r$-ball of radius
$R=1/(2k)$ about $\boldsymbol U=(I_d/k,\ldots,I_d/k)$.

\begin{theorem}[A dimension-independent interior comparison]
\label{thm:multioutcomemetric82}
For $d,N\ge1$, $k\ge2$, and
$\boldsymbol E,\boldsymbol F\in\mathfrak P_{d,k}$,
\begin{equation}\label{eq:multimetric82}
 \frac1{128}\min\{1,\sqrt N\,r(\boldsymbol E,\boldsymbol F)\}
 \le D_N^{\rm na}(\boldsymbol E,\boldsymbol F)
 \le D_N(\boldsymbol E,\boldsymbol F)
 \le\min\{2,2k\sqrt N\,r(\boldsymbol E,\boldsymbol F)\}.
\end{equation}
The lower bound holds for arbitrary legal ordered $k$-outcome
measurements, without the spectral bounds in
\eqref{eq:balancedbody82}.  The upper constants are independent of
$d$ and $N$; dependence on $k$ is displayed and is not asserted sharp.
\end{theorem}
\begin{proof}
Let $H_j=F_j-E_j$ and $A_j(t)=E_j+tH_j$ for $0\le t\le1$.
Each $A_j(t)$ is positive definite.  Solve the matrix differential
equation
\begin{equation}\label{eq:horizontalode82}
 \dot X_j(t)=\tfrac12 X_j(t)A_j(t)^{-1}H_j,
 \qquad X_j(0)=\sqrt{E_j}.
\end{equation}
Its coefficients are continuous.  Differentiating $X_j^*X_j$ shows
that it solves the linear equation
\[
 \dot B_j=\tfrac12 H_jA_j^{-1}B_j
             +\tfrac12 B_jA_j^{-1}H_j,\qquad B_j(0)=E_j.
\]
The function $A_j$ solves this same equation, so uniqueness gives
$X_j^*X_j=A_j$.  Consequently
\begin{equation}\label{eq:horizontalidentities82}
 X_j^*\dot X_j=\tfrac12 H_j,\qquad
 \dot X_j^*\dot X_j=\tfrac14 H_jA_j^{-1}H_j.
\end{equation}
This argument neither diagonalizes the tuple nor assumes that $A_j$
and $H_j$ commute.

Use the Stinespring isometry
\[
 W_t\psi=\sum_{j=1}^k |j\rangle_Y|j\rangle_Z X_j(t)\psi.
\]
Tracing out $Z$ and the last factor gives
$\mathcal M_{\boldsymbol A(t)}$.  The environment in this dilation
is used only in the proof and is not supplied by the device.  Since
$\sum_jH_j=0$, \eqref{eq:horizontalidentities82} yields
\begin{equation}\label{eq:horizontalzero82}
 W_t^*W_t=I_d,\qquad W_t^*\dot W_t=0,\qquad
 \dot W_t^*\dot W_t=\tfrac14\sum_jH_jA_j(t)^{-1}H_j.
\end{equation}

For completeness, purify a common adaptive tester, retaining a fresh
inaccessible environment for each call.  Its controls are isometries
independent of $t$.  Differentiate the final purified state.  It is a
sum of at most $N$ terms, with $\dot W_t$ replacing one occurrence
of $W_t$ in each term.  In the inner product of two distinct terms,
cancel the later common isometries.  At the later differentiated
slot the remaining contraction is $W_t^*\dot W_t$ or its adjoint,
hence zero.  Each diagonal term has norm at most
$\|\dot W_t\|_{\rm op}$.  The derivative of the purified state
therefore has norm at most $\sqrt N\|\dot W_t\|_{\rm op}$.
The derivative of its rank-one density operator has trace norm at
most twice this number.  Partial trace and all final processing
contract trace norm.  Bounded stopping is covered by a coherent
stopping flag and dummy calls on a fixed discarded input after
stopping; the resulting final state of the original tester is a
common postprocessing of this padded experiment.

Integrating the derivative bound gives the more informative inequality
\begin{align}\label{eq:multipath82}
 D_N(\boldsymbol E,\boldsymbol F)
 &\le \sqrt N\int_0^1
       \left\|\sum_jH_jA_j(t)^{-1}H_j\right\|_{\rm op}^{1/2}\,dt\\
 &\le\sqrt{2kN}\left\|\sum_jH_j^2\right\|_{\rm op}^{1/2}
 \le\sqrt2 k\sqrt N\max_j\|H_j\|_{\rm op}.\notag
\end{align}
We used $A_j(t)^{-1}\preceq2kI_d$ and the order-preserving map
$B\mapsto H_jBH_j$; commutation is unnecessary.  The weaker
constant $2k$ in \eqref{eq:multimetric82} is convenient below.
The trace distance is always at most two.

For the lower bound, choose a label $j$ and a unit eigenvector of
$F_j-E_j$ attaining its operator norm.  Send that same vector in
every call and replace the observed label by the indicator of $j$.
The two records are products of Bernoulli laws with parameters
$s,t\in[0,1]$ and $|s-t|=\|F_j-E_j\|_{\rm op}$.  Since
$\max\{s(1-s),t(1-t)\}+N^{-1}\le5/4$, the quantity in
\eqref{eq:spectralmodulus75} is at least
$\tfrac12\sqrt N|s-t|$.
Lemma~\ref{lem:bernoulliproduct75} gives the lower constant $1/128$.
Maximize over $j$.  No interior assumption was used in this part.
\end{proof}

\subsection{The affine norm ball and its entropy}
Let $\mathcal C_{N,k}(\delta)$ be the least number of radius-$\delta$
$D_N$-balls covering $\mathfrak P_{d,k}$.  Their centres may be
arbitrary legal memoryless ordered $k$-outcome measurements.

\begin{theorem}[Dimension-uniform finite-outcome entropy]
\label{thm:multioutcomeentropy82}
For $d,N\ge1$, $k\ge2$, and $0<\delta<1/128$,
\begin{equation}\label{eq:multientropy82}
 \max\left\{1,\left(\frac{\sqrt N}{256k\delta}\right)^p\right\}
 \le\mathcal C_{N,k}(\delta)
 \le\left(\frac{3\sqrt N}{\delta}\right)^p,
 \qquad p=(k-1)d^2.
\end{equation}
In particular, uniformly in $d,N,\delta$ for each fixed $k$,
\begin{equation}\label{eq:multientropylog82}
 \log_2\mathcal C_{N,k}(\delta)
 =\frac p2\log_2N+p\log_2(1/\delta)+O(p\log(k+1)).
\end{equation}
\end{theorem}
\begin{proof}
If $D_N(\boldsymbol E,\boldsymbol F)\le\delta<1/128$, the lower
bound of Theorem~\ref{thm:multioutcomemetric82} implies
$r(\boldsymbol E,\boldsymbol F)\le128\delta/\sqrt N$.
Lebesgue volume in the $p$-dimensional translated affine space gives
the lower covering bound, because the target is its $r$-ball of
radius $1/(2k)$.  This argument also applies when the centre lies
outside $\mathfrak P_{d,k}$.

Take a maximal $\epsilon$-separated subset of that ball, where
$\epsilon=\delta/(2k\sqrt N)$.  It is an $r$-net.  The open
$r$-balls of radius $\epsilon/2$ about its points are disjoint and
lie in the ball of radius $R+\epsilon/2$.  Their number is at most
$(1+2R/\epsilon)^p=(1+2\sqrt N/\delta)^p$.
By the upper comparison it is a $D_N$-cover at radius $\delta$.
This proves \eqref{eq:multientropy82}; taking logarithms gives
\eqref{eq:multientropylog82}, including the range where the displayed
power in the lower bound is less than one.
\end{proof}

\subsection{A finite grid with exact normalization}
An independent rounding of all $k$ effects need not preserve their
sum.  The following grid preserves both normalization and the two
spectral inequalities simultaneously.

\begin{lemma}[Rational nets of the balanced body]
\label{lem:multigrid82}
Let $0<t\le1/(8k)$ be rational and choose
\begin{equation}\label{eq:multigriddenominator82}
 K=\left\lceil\max\left\{\frac{4(k-1)d}{t},\,8k(k-1)d\right\}\right\rceil.
\end{equation}
Enumerate the first $k-1$ Hermitian matrices using diagonal
coordinates $0,1/K,\ldots,1$ and real and imaginary coordinates
above the diagonal in $K^{-1}\mathbb Z\cap[-1,1]$.
Set the last matrix equal to $I_d$ minus their sum and retain
exactly the legal tuples in $\mathfrak P_{d,k}$.  This finite
ordered grid, denoted $\mathcal G(t)$, is an $r$-net of radius $t$.
All membership and pairwise norm-threshold tests are finite exact
rational positive-semidefinite tests.
\end{lemma}
\begin{proof}
Put $e=(k-1)d/K$ and $\gamma=4ke\le1/2$.
For a target tuple form
$E'_j=(1-\gamma)E_j+\gamma I_d/k$.
Each $E'_j$ has margin $2e$ from both spectral faces in
\eqref{eq:balancedbody82}, and $\|E'_j-E_j\|_{\rm op}\le2e$.
Round the independent real coordinates of the first $k-1$ matrices
to nearest multiples of $1/K$, with a fixed tie rule and Hermitian
conjugacy.  The row-sum bound gives operator error at most $d/K$
for each rounded matrix.  Define the last matrix by exact
subtraction from $I_d$; its error is at most $(k-1)d/K=e$.
Thus every rounded effect still has margin at least $e$, their sum
is exactly $I_d$, and their distance from the target is at most
$3e\le3t/4<t$.  The rounded coordinates occur in the stated enumeration.

Legality is checked on
$E_j-I_d/(2k)$ and $3I_d/(2k)-E_j$, for every $j$.
The relation $r(\boldsymbol E,\boldsymbol F)\le s$ is decided by
checking $sI_d\pm(E_j-F_j)\succeq0$ for every $j$.
Exact Schur elimination, including the zero-pivot case, decides
these tests over the Gaussian rationals.  Equality is included.
\end{proof}

\begin{theorem}[A public exact code]
\label{thm:multioutcomecode82}
For integers $d,N\ge1$, $k\ge2$, and rational
$0<\delta<1/128$, there is a finite deterministic dictionary with
centres in $\mathfrak P_{d,k}$ and covering error at most $\delta/2$.
Its fixed-length index uses
\begin{equation}\label{eq:multicodebits82}
 B=\frac p2\log_2N+p\log_2(1/\delta)+O(p\log(k+1))
\end{equation}
bits, with absolute implied constants.  A fixed public decoder into
arbitrary legal $k$-outcome measurements cannot improve this order
for worst-case radius $\delta$.  The construction is finite and
exact; no polynomial running-time or workspace assertion is made.
\end{theorem}
\begin{proof}
Put $n=\lceil\sqrt N\rceil$ and $s=\delta/(8kn)$, and enumerate
$\mathcal G(s)$ from Lemma~\ref{lem:multigrid82}.
Retain a point precisely when its $r$-distance from every earlier
retained point is strictly greater than $s$.  The list is finite
and nonempty.  It is an $r$-cover of the grid at radius $s$;
together with the grid approximation it covers the body at radius
$2s$.  Its $D_N$-covering error is therefore at most
$2k\sqrt N(2s)\le\delta/2$.

Distinct retained points have $r$-distance greater than $s$.
The volume argument of Theorem~\ref{thm:multioutcomeentropy82}
gives dictionary size at most
\[
 (1+2R/s)^p=(1+8n/\delta)^p
             \le(17\sqrt N/\delta)^p.
\]
The lower bound follows from the same theorem because the dictionary
covers at radius $\delta$, even allowing arbitrary legal centres.
Taking the ceiling of the logarithm changes the length by at most one.

Encoder and decoder reconstruct the same ordered dictionary from
$(d,k,N,\delta)$; only its index is transmitted.  Unused binary
words decode to the fixed uniform tuple.  For a supplied rational
target, scan the list for the first centre within $2s$, using the
exact tests in Lemma~\ref{lem:multigrid82}.
For a real target supplied with certified rational approximations,
search for a grid point at certified strict distance less than $s$
and use its assigned predecessor.  Such a point exists at distance
at most $3s/4$ by the grid lemma, so progressively finer certified
input approximations make the strict test terminate.  An unrestricted oracle for arbitrary real
coordinates is not part of the encoding model.  The covering
assertion itself concerns every real target, irrespective of its
computational representation.
\end{proof}

The grid has at most
$[(K+1)^d(2K+1)^{d(d-1)}]^{k-1}$ candidates before the legality
tests.  Reconstructing this list and comparing its retained points
can cost far more than the transmitted index.  The estimate concerns
the description of an unknown measurement, not free communication
of a target-dependent dictionary.

\section{Joint learning and finite certificates for a finite alphabet}
\label{sec:multilearning82}

We next construct one learner common to all unknown elements of
$\mathfrak P_{d,k}$.  Its training experiments do not use the
pair-dependent eigenvectors in the geometric lower bound.
The outcome count $k$ is public, as are $d,N,\delta,\eta$.
The upper procedure uses independent one-call probe--reference pairs;
coherent collective processing is confined to completed outputs.
The converse permits arbitrary coherent adaptive training, with an
every-record budget and bounded public stopping.

\begin{theorem}[Joint finite-outcome learning and coding]
\label{thm:multioutcomelearning82}
Let $d,N\ge1$, $k\ge2$, $0<\delta\le2^{-24}$, and
$0<\eta\le1/8$.  Write $M^*_{d,k}(N,\delta,\eta)$ for the
minimum every-record call budget with worst-case failure at most
$\eta$ under $D_N$-error $\delta$ on $\mathfrak P_{d,k}$.
There are absolute constants $c,C>0$ such that
\begin{equation}\label{eq:multilearning82}
 c\frac{N}{\delta^2}\left(d^2+d\log\frac1\eta\right)
 \le M^*_{d,k}(N,\delta,\eta)
 \le C\frac{k^3N}{\delta^2}
                  \left(d^2+d\log\frac{k}{\eta}\right).
\end{equation}
For each fixed $k$, this is a matching order jointly in
$d,N,\delta,\eta$, including $d=1$:
\[
 M^*_{d,k}(N,\delta,\eta)
 =\Theta_k\!\left(N\delta^{-2}
                   [d^2+d\log(1/\eta)]\right).
\]
The displayed dependence on growing $k$ is not claimed optimal.

For rational accuracy and confidence, the upper procedure has a
finite deterministic specification, returns an index with a fixed
public decoder into $\mathfrak P_{d,k}$, and simultaneously uses
\begin{equation}\label{eq:multilearnedbits82}
 B=\frac{(k-1)d^2}{2}\log_2N
       +(k-1)d^2\log_2(1/\delta)+O((k-1)d^2\log(k+1))
\end{equation}
bits.  On its good event the decoded error is at most $5\delta/8$.
The same call order is available with finite trusted controls.
Classical search, readout storage, known-state preparation, and
known-control realization are separate resources.
\end{theorem}
\begin{proof}
We first give the constructive upper bound for rational parameters.
Set $n=\lceil\sqrt N\rceil$ and
\begin{equation}\label{eq:multitrainingaccuracy82}
 \epsilon=\frac{\delta}{16kn}.
\end{equation}
For each $j$ consider the binary effect
$F_j=I_d/4+E_j/2$.  It belongs to $[I_d/4,3I_d/4]$.
One call to $\mathcal M_{\boldsymbol E}$ simulates a call to
$\mathcal M_{F_j}$: conditional on the original label $x$, return
bit one with probability $3/4$ when $x=j$ and with probability
$1/4$ otherwise.  This is an equality of channels, including their
action on a retained reference, because the bit-one effect is
\[
 \tfrac34E_j+\tfrac14\sum_{x\ne j}E_x=\tfrac14I_d+\tfrac12E_j.
\]
Only public classical randomness is added.  Each simulated call
uses exactly one call to the unknown device.

Apply Theorem~\ref{thm:rationallearner81} to each $F_j$ at future
horizon one, accuracy $\epsilon$, and failure $\eta/k$, using
fresh records for different $j$.  Its output $\widehat F_j$ is
Gaussian rational and legal.  The binary identity
$d_1(F_j,\widehat F_j)=2\|F_j-\widehat F_j\|_{\rm op}$ implies
that, with probability at least $1-\eta$, all the Hermitian
matrices $B_j=2\widehat F_j-I_d/2$ obey
\begin{equation}\label{eq:multiassembly82}
 \max_j\|B_j-E_j\|_{\rm op}\le\epsilon.
\end{equation}
Only a union bound is needed; no independence assertion concerning
selected quantum readout outcomes is used.

The matrices $B_j$ need not sum to $I_d$.  Reconstruct the finite
legal grid $\mathcal G(\epsilon/2)$ and choose its first tuple
$\boldsymbol G$ satisfying
\[
 \|G_j-B_j\|_{\rm op}\le2\epsilon\quad\text{for every }j.
\]
Each comparison is an exact rational positive-semidefinite test.
On \eqref{eq:multiassembly82}, such a tuple exists by
Lemma~\ref{lem:multigrid82}; it has
$r(\boldsymbol G,\boldsymbol E)\le3\epsilon$.
On every other record, if the finite search finds no tuple, return
$\boldsymbol U$.  Thus every record has a legal output and the
same predetermined training budget.  The geometric upper bound gives
\[
 D_N(\boldsymbol E,\boldsymbol G)
       \le6k\sqrt N\epsilon\le3\delta/8.
\]
Encode $\boldsymbol G$ with the dictionary of
Theorem~\ref{thm:multioutcomecode82} at parameter $\delta/2$.
Its additional error is at most $\delta/4$, so the decoded error
is at most $5\delta/8$.  It uses no further device calls and has
length \eqref{eq:multilearnedbits82}.

There are $k$ binary learners, each using at most
$C_0\epsilon^{-2}[d^2+d\log(k/\eta)]$ calls.  Since $n^2\le4N$,
this proves the upper bound.  The inherited rational synthesis
and finite-control theorem applies to the simulated binary channels
as to any channels of the same form.  The public $1/4,3/4$ relabeling
is exact classical control.  The error guarantee used in
\eqref{eq:multiassembly82} is the inherited actual finite-control
guarantee, not an idealized readout subsequently assumed exact.
The construction terminates in finite classical time, with no
assertion that this time is polynomial.  For nonrational error
parameters, choose rational $\delta'\in[\delta/2,\delta]$ and
$\eta'\in[\eta/2,\eta]$ to obtain the stated existence bound;
these choices change only absolute constants.

For the lower bound let $m=\lfloor k/2\rfloor$ and $q=2m/k\ge2/3$.
Embed any binary interior effect $F\in[I_d/4,3I_d/4]$ by the tuple
\begin{equation}\label{eq:multiembedding82}
 E_j(F)=\begin{cases}
             2F/k,&1\le j\le m,\\
             2(I_d-F)/k,&m<j\le2m,\\
             I_d/k,&j=k=2m+1.
             \end{cases}
\end{equation}
This belongs to $\mathfrak P_{d,k}$.  A query to this tuple is
simulated using one binary $F$ query: with probability $q$, retain
the binary outcome and select a uniform label from its corresponding
group of $m$ labels; otherwise erase the bit and return the extra
label.  In the odd case the erased reference state is the sum of
its two subnormalized binary states, as required by the effect
$I_d/k$.  Thus the simulation is valid for entangled inputs and
adaptive controls, not merely for outcome probabilities on pure
states.  For even $k$ there is no erasure.

Suppose a $k$-outcome learner returns any legal tuple
$\widehat{\boldsymbol E}$ within $D_N$-error $\delta$.
Coarse-graining the first $m$ labels gives a binary effect
$H=\sum_{j\le m}\widehat E_j$ satisfying
$d_N(qF,H)\le\delta$.  This follows by executing the same
coarse-graining inside any binary tester of the two induced channels.
The Bernoulli lower bound in
Theorem~\ref{thm:multioutcomemetric82}, applied to these two binary
effects, implies
\[
 \|H-qF\|_{\rm op}\le128\delta/\sqrt N,
\]
since $\delta<1/128$.
Spectrally clip $H/q$ to $[I_d/4,3I_d/4]$, obtaining $\widehat F$.
Weyl's inequality bounds the clipping displacement by
$\|H/q-F\|_{\rm op}$; the triangle inequality therefore gives
\[
 \|\widehat F-F\|_{\rm op}
      \le\frac{256}{q}\frac\delta{\sqrt N}
      \le\frac{384\delta}{\sqrt N}.
\]
No operator-Lipschitz claim about spectral clipping is used.
The last accuracy is at most $2^{-14}$ under the stated cap.
The proof of Corollary~\ref{cor:operatorconfidence80} gives the
binary operator-loss lower bound already on this interior, with
arbitrary coherent adaptive training.  Applying it to the simulation
proves the left inequality of \eqref{eq:multilearning82}, including
the contribution $d\log(1/\eta)$ and the scalar case.  The same
simulation preserves every-record budgets and public stopping.
\end{proof}

The length in \eqref{eq:multilearnedbits82} has the optimal leading
order also among learners with a fixed-length, fixed public decoder.
Indeed, if failure is less than one at every target, at least one
word must decode inside that target's loss ball.  The decoder's
image is therefore a covering of the entire target family.
The lower bound in Theorem~\ref{thm:multioutcomeentropy82} applies
regardless of the number of training calls.  This is separate from
the statistical query lower bound.  Shared-randomness expected-prefix
conventions are different and are not part of this assertion.

\subsection{Finite uniform-risk certificates}
The certification mechanism also survives a finite outcome alphabet.
For an ideal one-call maximally entangled acquisition, its normalized
classical--quantum record is
\begin{equation}\label{eq:multichoi82}
 \Omega(\boldsymbol E)=\frac1d\bigoplus_{j=1}^k E_j^{\mathsf T}.
\end{equation}
The classical labels remain ordered.  A collective readout after $M$
independent acquisitions can be described by one POVM
$(A_{c,x})_c$ on dimension $d^M$ for each classical string
$x\in\{1,\ldots,k\}^M$.  Here $A_{c,x}\succeq0$ and
$\sum_c A_{c,x}=I_{d^M}$.  Its output label $c$ is decoded by a
fixed finite family $\boldsymbol C_c$ of legal $k$-outcome tuples.

\begin{theorem}[Finite-alphabet certificate transfer]
\label{thm:multicertificate82}
Let $\mathcal G\subset\mathfrak P_{d,k}$ be a complete finite
$r_0$-net in the tuple norm $r$.  Suppose a fixed readout satisfies
\begin{equation}\label{eq:multicertificate82}
 \sum_{c:\,r(\boldsymbol G,\boldsymbol C_c)\le a}
 \sum_{x\in\{1,\ldots,k\}^M}
 \operatorname{tr}\left[A_{c,x}
        d^{-M}\bigotimes_{i=1}^M G_{x_i}^{\mathsf T}\right]
 \ge1-\alpha\qquad(\boldsymbol G\in\mathcal G).
\end{equation}
Then every real $\boldsymbol E\in\mathfrak P_{d,k}$ satisfies
\begin{equation}\label{eq:multitransfer82}
 \mathbb P_{\boldsymbol E}
   \{r(\boldsymbol E,\boldsymbol C_J)>a+r_0\}
 \le\alpha+\frac{Mk r_0}{2}.
\end{equation}
If all decoding centres are in $\mathfrak P_{d,k}$, the same failure
bound holds for $D_N$-error greater than $2k\sqrt N(a+r_0)$.
For Gaussian-rational data, rational $a,r_0,\alpha$, and the net
of Lemma~\ref{lem:multigrid82}, the hypotheses admit a finite exact
verification.  A proper prefix of the net or of the classical
string list is not a certificate.
\end{theorem}
\begin{proof}
The block-diagonal formula \eqref{eq:multichoi82} gives
\[
 \|\Omega(\boldsymbol E)-\Omega(\boldsymbol G)\|_1
    =d^{-1}\sum_j\|E_j-G_j\|_1
    \le k\,r(\boldsymbol E,\boldsymbol G).
\]
Tensor telescoping between normalized states gives
\begin{equation}\label{eq:multirecordcontinuity82}
 \|\Omega(\boldsymbol E)^{\otimes M}
           -\Omega(\boldsymbol G)^{\otimes M}\|_1
       \le Mk\,r(\boldsymbol E,\boldsymbol G).
\end{equation}
Apply the same readout to both states.  Its index laws differ in
total variation by at most $Mk r_0/2$ when
$r(\boldsymbol E,\boldsymbol G)\le r_0$.
Choose such a $\boldsymbol G$.  The fixed set of labels
$\{c:r(\boldsymbol G,\boldsymbol C_c)\le a\}$ has probability
at least $1-\alpha$ at $\boldsymbol G$, by
\eqref{eq:multicertificate82}; every label in this set has
$r(\boldsymbol E,\boldsymbol C_c)\le a+r_0$.
Transferring this one event proves \eqref{eq:multitransfer82}.
There is no assumption that the success-label sets vary continuously
with the unknown target.

The geometric upper comparison proves the assertion about $D_N$.
For exact verification, reconstruct the whole legal net and every
classical string.  All traces in \eqref{eq:multicertificate82} are
Gaussian-rational expressions with rational real values.  Check
complete POVM positivity and normalization at every string, use
$sI_d\pm(G_j-C_{c,j})\succeq0$ to determine the good labels
including equality, and check each rational inequality.  These are
finite computations.  Their soundness on the real parameter body
is supplied by \eqref{eq:multirecordcontinuity82}, not by a numerical
sampling argument.
\end{proof}

The certificate describes the ideal specified rational readout.
When implemented controls have a separately certified output-law
error $\zeta$, the right side of \eqref{eq:multitransfer82} increases
by at most $\zeta$.  Such an error must be budgeted, not silently
identified with zero.  The finite-control guarantee in
Theorem~\ref{thm:multioutcomelearning82} instead uses the already
budgeted binary procedures.  The exact certificate transfer and
the actual-control learning guarantee thus have explicitly distinct
hypotheses.

\section{Search-free legalization of balanced estimates}
\label{sec:affinerepair83}

The exact public dictionary and the statistical legalization step
have different purposes.  A dictionary exploits global packing to
reduce transmitted length.  Legalizing a collection of component
estimates need not reconstruct such a dictionary.

\begin{theorem}[Affine repair with a uniform error bound]
\label{thm:affinerepair83}
Let $B=(B_1,\ldots,B_k)$ be Hermitian matrices and $a>0$.  Define
\begin{align}
 A_j&=B_j+\frac{I_d-\sum_iB_i}{k},\label{eq:affineprojection83}\\
 L_a(B)_j&=\frac{A_j+4a I_d}{1+4ka}.\label{eq:affinerepair83}
\end{align}
Suppose $E\in\mathfrak P_{d,k}$, the balanced interior from
\eqref{eq:balancedbody82}, and
$\max_j\|B_j-E_j\|_{\rm op}\le a$.
Then $L_a(B)\in\mathfrak P_{d,k}$ and
\begin{equation}\label{eq:affineerror83}
 \max_j\|L_a(B)_j-E_j\|_{\rm op}
 \le\frac{4a}{1+4ka}\le4a.
\end{equation}
For rational input and rational $a$, this is a rational map.  Testing
its balanced legality and returning the uniform tuple upon failure
defines a total legal procedure on every Hermitian input record.
It uses $O(kd^2)$ rational arithmetic operations apart from $2k$
exact positive-semidefinite tests of order $d$, and polynomial bit
complexity.  There is no grid enumeration.
\end{theorem}
\begin{proof}
Let $\Delta_j=B_j-E_j$.  Then
$A_j-E_j=\Delta_j-k^{-1}\sum_i\Delta_i$, whose norm is at most $2a$,
and $\sum_jA_j=I_d$.  Consequently
\[
 (1/(2k)-2a)I_d\preceq A_j\preceq(3/(2k)+2a)I_d.
\]
Adding $4aI_d$ and dividing by $1+4ka$ gives exactly the two
balanced spectral inequalities.  The sum of the repaired effects
is $I_d$.  Finally,
\[
 L_a(B)_j-E_j
 =\frac{A_j-E_j+4ka(I_d/k-E_j)}{1+4ka}.
\]
The balanced hypothesis gives $\|I_d/k-E_j\|_{\rm op}\le1/(2k)$,
proving \eqref{eq:affineerror83}.  Rationality and the operation count
follow from the formula.  The exact legality test makes no good-event
assumption; if it fails the uniform tuple is an explicitly specified
legal output.  On the good event it never fails.
\end{proof}

\begin{corollary}[A common learner with polynomial legalization]
\label{cor:affinelearner83}
Under the range and access conventions of
Theorem~\ref{thm:multioutcomelearning82}, replace its tuple-grid
legalization by the guarded map in Theorem~\ref{thm:affinerepair83}.
Taking
\[
 a=\frac{\delta}{32k\lceil\sqrt N\rceil}
\]
and allocating failure $\eta/k$ to each fresh binary component
learner gives, on every record,
\[
 M\le Ck^3N\delta^{-2}[d^2+d\log(k/\eta)]
\]
for an absolute $C$, with future-use error at most $\delta/4$ on a
common event of probability at least $1-\eta$.
Encoding with the public dictionary at parameter $\delta/2$ gives
error at most $\delta/2$ on the same event, without further calls,
and retains the optimal fixed-$k$ fixed-length description order.
Only the legalization cost is made polynomial; neither dictionary
construction nor general collective-readout synthesis is asserted
to be efficient.
\end{corollary}
\begin{proof}
Use precisely the reference-preserving binary simulation and actual
finite-control component learners from
Theorem~\ref{thm:multioutcomelearning82}, at the displayed smaller
accuracy.  Its conditional good-event argument gives component
error at most $a$ with total failure at most $\eta$.  Affine repair
then gives tuple error at most $4a$.  The balanced metric upper
bounds the future-use error by
$8k\sqrt N a\le\delta/4$.  The dictionary adds at most $\delta/4$.
Fresh calls and all finite-control budgets are charged exactly as in
the component theorem; decreasing the component accuracy by a factor
two changes its call bound only by an absolute factor.  The guarded
map is total even off the good event, so the recordwise call bound
and public stopping convention are unchanged.
\end{proof}

\section{Relation to channel metrology and measurement information}
\label{sec:comparison84}

\subsection{Support spans, channel estimation and finite pairs}
For the measurement Kraus operators
$L_{j,a}=|j\rangle\langle a|\sqrt{E_j}$, the real Hermitian
Kraus-product span is exactly $\mathcal W_E$.  For the orbit
$L_{j,a}(t)=L_{j,a}e^{-itG}$ the effective generator is $G$.
Zhou--Jiang \cite[Theorems~1--3 and Section~V]{ZJ84} establish the
Kraus-span criterion for the linear or quadratic asymptotic channel
Fisher-information rate and give error-correction attainability.
We do not claim a new metrological dichotomy or a new correction
principle.  Theorems~\ref{thm:supportkernel84} and
\ref{thm:orbittrichotomy84} identify its complete support-kernel
form and give finite-angle trace-distance inequalities with an
explicit $4mt^2\|G\|_{\mathrm{op}}^2$ remainder.
Theorem~\ref{thm:curveclassification85} adds a normalized first-order
realization for general two-sided $C^2$ effect curves; its
$\sqrt N|t|+O(Nt^2)$ bound handles rank changes without a smooth
exact Kraus hypothesis.  Theorem~\ref{thm:controlstability85} separately
budgets certified control errors.  Neither constitutes a new claim
to the underlying metrological exponents.

Sieniawski--Demkowicz-Dobrza\'nski
\cite[Section~IV, equations~(16)--(23)]{SD76} integrate channel
Fisher-information bounds to bound adaptive discrimination and
compare metrological scaling with perfect finite discrimination.
The same horizontal method is a principal antecedent of
\eqref{eq:frontupper84}.  Our support-kernel identity is stated in
measurement coordinates and our lower \eqref{eq:finiteanglelower84}
comes from an actual corrected finite-angle channel.  It does not
claim the exact optimal tester or perfect discrimination.

\subsection{Fisher frames and covariance embeddings}
Saini--Kiukas--Burgarth--Gilchrist \cite[Theorem~1]{SKBG84} compare
classical and quantum Fisher information for a varying full-rank
\emph{state} measured by one fixed informationally complete POVM,
using a state-dependent frame operator.  Here the varying object
is the measurement, the tangent has coupled normalization, and the
loss permits $N$ adaptive uses with a reference.  In the rank-one
case our support span equals the effect span, so informational
completeness excludes the linear orbit regime by
Corollary~\ref{cor:orbitexamples84}.  That intersection does not
identify the two quadratic forms or their resource models.

Mayer--Yun \cite[Theorem~3.6, Definition~3.7 and Remark~4.4]{MY84}
use a kernel embedding of operator-valued measures, a kernel-dependent
discrepancy, and a tomographic Gram operator whose nullspace records
unobserved state directions.  Our support span contains entire
matrix blocks $P_j\mathcal H_dP_j$, rather than one observable
per effect; this distinction disappears for rank-one effects but
not in general.  Theorem~\ref{thm:supportkernel84} concerns a
coupled measurement tangent, and \eqref{eq:frontupper84} concerns
adaptive finite-use loss.  Neither kernel embedding nor injectivity
alone supplies those statements.

\subsection{Known symmetry and tomography primitives}
The full primary text of Yoshida--Okigami--Posta--Grinko
\cite{YOPG84} distinguishes known-symmetry estimation from known-pair
discrimination.  Theorem~7 gives invariant-state copy complexity
$\Theta((D_{\rm st}-1+\log(1/\eta))/\epsilon^2)$, where
$D_{\rm st}=\sum_\lambda m_\lambda^2$; $D_{\rm st}=1$ is a
known singleton.  Corollary~9 gives parallel covariant-channel bounds
\[
 O((D'_G+K_G\log(1/\eta))/\epsilon_\diamond^2),\qquad
 O((C'_G+\log(1/\eta))/\epsilon_{\rm tr}^2),
\]
in one-use halved diamond distance and normalized-Choi trace distance,
respectively, for $0<\epsilon\le1$, $0<\eta\le1/2$.
Here $D_G=\sum_\lambda m_\lambda M_\lambda$ counts Choi commutant
parameters before trace preservation,
$K_G=\sum_\lambda m_\lambda$,
$D'_G=D_G+K_G\log(4L)$, and
$C'_G=d_{\mathcal I}^{-1}(\sum_\lambda m_\lambda
\sqrt{d_\lambda M_\lambda})^2$; $L$ counts occupied blocks.
Theorem~13 does not give a general matching inverse-square diamond
lower: its bound is $\Omega_\xi(D_G/\epsilon_\diamond^{\xi/2})$
under its affine-dimension hypothesis.

For permutation covariance on $\ell$ systems of fixed local dimension
$q$, Theorems~12 and~15 match the Choi order
$\ell^{q^4-q^2}\epsilon_{\rm tr}^{-2}$ at constant confidence.
Theorem~18 approximates the Hayashi measurement in
$O((s+q)\operatorname{polylog}(s,q,\tau^{-1},\zeta^{-1}))$ gates,
with coupled output error at most $\tau$ except with probability
$\zeta$.  This is not a diamond-error certificate for our recovery.
Known-symmetry estimation, its parameter dimensions, and that
specialized gate construction do not themselves give our support
kernel or finite-pair curve theorem.  Nor do our theorems claim
faster learning on those symmetry classes.

For ordinary tomography the main-text substitutions remain explicit.
Mele--Bittel \cite[Theorem~III.3]{MB67} give the sufficient
constant-confidence order $O(dkr\epsilon^{-2})$ in halved diamond
loss, with $r\le dk$ for measurement channels.
Zambrano--Ramos-Calderer--Kueng \cite[Theorem~2]{ZRK76} give the
sufficient single-copy worst-input outcome-TV bound
\[
 m\ge\frac{8[d^2+\epsilon(d^2+1)/6]}{\epsilon^2}
          \log\frac{2^{k+1}9^{2d}}{\eta}.
\]
Either loss bounds the component norm error by $\epsilon$ after
output dephasing where needed.  A legal estimator need not be
balanced: apply Theorem~\ref{thm:affinerepair83} before using the
interior metric.  The resulting future loss is at most
$8k\sqrt N\epsilon$.  With $\epsilon=\delta/(8k\sqrt N)$ the
respective sufficient orders are
\[
 O(d^2k^4N\delta^{-2}),\qquad
 O\!\left(k^2N\delta^{-2}d^2[d+k+\log(1/\eta)]\right).
\]
These are upper-bound substitutions, not optimality claims about
the cited estimators.  Our component procedure credits the existing
binary operator-norm primitive and uses a different confidence
allocation.  The full calculations and confidence restrictions
are retained in the supplement, Section~\ref{sec:priority83}.

\subsection{What remains beyond the theorems}
The orbit theorem and Theorem~\ref{thm:curveclassification85} provide
matching scales on each fixed two-sided smooth curve with nonzero
first derivative, including second-order rank openings.  Their
constants need not be uniform across curves.  Arbitrary-pair
midpoint equivalence and full-body multi-outcome entropy are not
consequences of these local statements.  Growing-$k$ minimax sharpness,
polynomial entropy-optimal dictionaries, and general readout synthesis
remain separate mathematical questions.  Independent expert priority
assessment is not inferred from the absence of an identical statement
in this targeted comparison.

\subsection{Tangent neighborhoods and independent acquisition}
Theorem~\ref{thm:tubedichotomy86} fixes $E,H$ and an explicit
second-order remainder allowance, rather than assuming a particular
smooth family of Kraus operators.  It treats the full one-sided
positive tangent cone.  A nonzero missing-support compression has a
classical impossible-event witness; only within the cone's lineality
space does the support-generator test select the coherent branch.
The latter still uses the credited correction mechanism.
Theorem~\ref{thm:producttrichotomy86} separates these mechanisms by
an explicit product probe--reference restriction.  Its upper uses
standard root-fidelity inequalities and tensor multiplicativity
\cite[Chapter~3]{Watrous65}, not a new fidelity principle or a bound
on all nonadaptive entangled strategies.

The higher-order statement requires the finite remainder
\eqref{eq:powerjet86}.  Proposition~\ref{prop:mixedrates86} instead
proves a two-scale law when an intermediate support-opening term is
present.  These results address one-sided openings and selected
higher-order approaches without identifying a first jet with an
arbitrary-pair global geometry.  No general boundary entropy or
unrestricted growing-outcome learning law is inferred.

\begin{thebibliography}{99}
\bibitem{Watrous65} J. Watrous, \emph{The Theory of Quantum Information},
Cambridge University Press, 2018. Theorems~2.22 and~2.26
(Choi criteria), Corollary~3.40 (trace-norm contraction),
and Section~3.3, equation~(3.306) (hybrid comparison).
\bibitem{PLS66} T. Surawy-Stepney, J. Kahn, R. Kueng and M. Gu\c{t}\u{a},
Projected least-squares quantum process tomography,
\emph{Quantum} \textbf{6} (2022), 844;
arXiv:2107.01060v2, 18 October 2022.
\bibitem{Gutoski66} G. Gutoski,
On a measure of distance for quantum strategies,
\emph{J. Math. Phys.} \textbf{53} (2012), 032202;
arXiv:1008.4636v4, 14 March 2012.
\bibitem{NZGB67} S. G. Naik, M. Zartab, N. Gisin and M. Banik,
No-Go Theorem for Generic Simulation of Qubit Channels with Finite
Classical Resources, arXiv:2501.15807v2 (16 July 2025),
Theorems~2 and~5.
\bibitem{MB67} A. A. Mele and L. Bittel,
Optimal learning of quantum channels in diamond distance,
arXiv:2512.10214v3 (2026), Theorem~I.1, Lemmas~III.7--III.8,
Corollary~III.9 and Sections~III.1.1, IV.1.1.
\bibitem{OG67} A. Oufkir and F. Girardi,
Improved Lower Bounds for Learning Quantum Channels in Diamond
Distance, arXiv:2601.04180v4 (2026).
\bibitem{DKG67} R. Demkowicz-Dobrza\'nski, J. Ko\l ody\'nski
and M. Gu\c{t}\u{a}, The elusive Heisenberg limit in quantum-enhanced
metrology, \emph{Nature Communications} \textbf{3} (2012), 1063.
DOI: 10.1038/ncomms2067; arXiv:1201.3940v2, equations~(4)--(6)
(classical simulation), equations~(10)--(13) (channel extension),
and the classical-simulation argument in Methods.
\bibitem{YBCH68} Y.~Yang, G.~Bai, G.~Chiribella and M.~Hayashi,
Compression for quantum population coding,
\emph{IEEE Transactions on Information Theory} \textbf{64} (2018), no.~7,
doi:10.1109/TIT.2017.2788407;
arXiv:1701.03372v8.
\bibitem{SHW68} F.~Salek, M.~Hayashi and A.~Winter,
Usefulness of adaptive strategies in asymptotic quantum channel
 discrimination, \emph{Phys. Rev. A} \textbf{105} (2022), 022419;
arXiv:2011.06569v3 (17 March 2024).
\bibitem{CMW70} T.~Cooney, M.~Mosonyi and M.~M.~Wilde,
Strong converse exponents for a quantum channel discrimination problem
and quantum-feedback-assisted communication,
arXiv:1408.3373v3 (26 February 2016), p.~5 and Sections~4.1--4.2.
\bibitem{Acin70} A.~Ac\'in,
Statistical distinguishability between unitary operations,
\emph{Phys. Rev. Lett.} \textbf{87} (2001), 177901.
DOI: 10.1103/PhysRevLett.87.177901.
\bibitem{DFY70} R.~Duan, Y.~Feng and M.~Ying,
Entanglement is not necessary for perfect discrimination between
unitary operations, \emph{Phys. Rev. Lett.} \textbf{98} (2007), 100503.
DOI: 10.1103/PhysRevLett.98.100503.

\bibitem{DW72} S. Das and M. M. Wilde,
Quantum reading capacity: General definition and bounds,
\emph{IEEE Trans. Inform. Theory} \textbf{65} (2019), 7566--7583;
arXiv:1703.03706v3; doi:10.1109/TIT.2019.2929925.
\bibitem{WBHK72} M. M. Wilde, M. Berta, C. Hirche and E. Kaur,
Amortized channel divergence for asymptotic quantum channel discrimination,
\emph{Lett. Math. Phys.} \textbf{110} (2020), 2277--2336;
arXiv:1808.01498v2; doi:10.1007/s11005-020-01297-7.
\bibitem{WW72} X. Wang and M. M. Wilde,
Resource theory of asymmetric distinguishability for quantum channels,
\emph{Phys. Rev. Research} \textbf{1} (2019), 033169;
arXiv:1907.06306v2; doi:10.1103/PhysRevResearch.1.033169.
\bibitem{PPKKO72} Z. Pucha\l a, \L. Pawela, A. Krawiec, R. Kukulski and
M. Oszmaniec, Multiple-shot and unambiguous discrimination of von Neumann
measurements, \emph{Quantum} \textbf{5} (2021), 425;
arXiv:1810.05122v4; doi:10.22331/q-2021-04-06-425.


\bibitem{SZ74} M. Sedl\'ak and M. Ziman,
Optimal single-shot strategies for discrimination of quantum measurements,
\emph{Phys. Rev. A} \textbf{90} (2014), 052312;
arXiv:1408.0934; doi:10.1103/PhysRevA.90.052312.
Section VI, Example 4, equation (37).
\bibitem{PPKK74} Z. Pucha\l a, \L. Pawela, A. Krawiec and R. Kukulski,
Strategies for optimal single-shot discrimination of quantum measurements,
\emph{Phys. Rev. A} \textbf{98} (2018), 042103;
arXiv:1804.05856; doi:10.1103/PhysRevA.98.042103. Theorem 1.
\bibitem{FM75} J.~Fiur\'a\v sek and M.~Mi\v cuda,
Optimal two-copy discrimination of quantum measurements,
\emph{Phys. Rev. A} \textbf{80} (2009), 042312;
arXiv:0909.2940; doi:10.1103/PhysRevA.80.042312.

\bibitem{KPP76} A.~Krawiec, \L.~Pawela and Z.~Pucha\l a,
Discrimination of POVMs with rank-one effects,
\emph{Quantum Information Processing} \textbf{19} (2020), 428;
arXiv:2002.05452v1; doi:10.1007/s11128-020-02883-3.
\bibitem{DBSA76} C.~Datta, T.~Biswas, D.~Saha and R.~Augusiak,
Perfect discrimination of quantum measurements using entangled systems,
\emph{New Journal of Physics} \textbf{23} (2021), 043021;
arXiv:2012.07069v2; doi:10.1088/1367-2630/abecaf.
\bibitem{FI76} A.~Fujiwara and H.~Imai,
A fibre bundle over manifolds of quantum channels and its application
to quantum statistics,
\emph{J. Phys. A: Math. Theor.} \textbf{41} (2008), 255304;
doi:10.1088/1751-8113/41/25/255304. Theorems~4 and~5.
\bibitem{DM76} R.~Demkowicz-Dobrza\'nski and L.~Maccone,
Using entanglement against noise in quantum metrology,
\emph{Phys. Rev. Lett.} \textbf{113} (2014), 250801;
arXiv:1407.2934v3; doi:10.1103/PhysRevLett.113.250801.
\bibitem{KGAD76} S.~Kurdzia\l ek, W.~G\'orecki, F.~Albarelli
and R.~Demkowicz-Dobrza\'nski,
Using adaptiveness and causal superpositions against noise in quantum metrology,
\emph{Phys. Rev. Lett.} \textbf{131} (2023), 090801;
arXiv:2212.08106v2; doi:10.1103/PhysRevLett.131.090801.
\bibitem{YF76} H.~Yuan and C.-H.~F.~Fung,
Fidelity and Fisher information on quantum channels,
arXiv:1506.00819v3 (2017), Section~III, equation~(26).
\bibitem{SD76} S.~Sieniawski and R.~Demkowicz-Dobrza\'nski,
Adaptive quantum channel discrimination using methods of quantum metrology,
\emph{New Journal of Physics} \textbf{28} (2026), 024502;
arXiv:2510.15506v3; doi:10.1088/1367-2630/ae3edd.
\bibitem{Dittmann76} J.~Dittmann,
Explicit formulae for the Bures metric,
\emph{J. Phys. A: Math. Gen.} \textbf{32} (1999), 2663--2670;
doi:10.1088/0305-4470/32/14/007;
arXiv:quant-ph/9808044, entitled
\emph{Note on explicit formulae for the Bures metric}.
\bibitem{SZGeom76} H.-J.~Sommers and K.~\.{Z}yczkowski,
Bures volume of the set of mixed quantum states,
\emph{J. Phys. A: Math. Gen.} \textbf{36} (2003), 10083--10100;
arXiv:quant-ph/0304041v3; doi:10.1088/0305-4470/36/39/308.
\bibitem{ZRK76} L.~Zambrano, S.~Ramos-Calderer and R.~Kueng,
Fast quantum measurement tomography with optimal error bounds,
\emph{Quantum} \textbf{10} (2026), 2162;
arXiv:2507.04500v3; doi:10.22331/q-2026-07-15-2162.
\bibitem{BDPS76} A.~Bisio, G.~M.~D'Ariano, P.~Perinotti and M.~Sedl\'ak,
Quantum learning algorithms for quantum measurements,
\emph{Physics Letters A} \textbf{375} (2011), 3425--3434;
arXiv:1103.0480v2; doi:10.1016/j.physleta.2011.08.002.
\bibitem{HB76} B.~L.~Higgins, D.~W.~Berry, S.~D.~Bartlett,
M.~W.~Mitchell, H.~M.~Wiseman and G.~J.~Pryde,
Demonstrating Heisenberg-limited unambiguous phase estimation
without adaptive measurements,
\emph{New Journal of Physics} \textbf{11} (2009), 073023;
arXiv:0809.3308v4; doi:10.1088/1367-2630/11/7/073023.
\bibitem{KLY76} S.~Kimmel, G.~H.~Low and T.~J.~Yoder,
Robust calibration of a universal single-qubit gate set via robust phase estimation,
\emph{Phys. Rev. A} \textbf{92} (2015), 062315;
doi:10.1103/PhysRevA.92.062315;
erratum, \emph{Phys. Rev. A} \textbf{104} (2021), 069901,
doi:10.1103/PhysRevA.104.069901; arXiv:1502.02677v3.

\bibitem{Hyllus78} P.~Hyllus, W.~Laskowski, R.~Krischek,
C.~Schwemmer, W.~Wieczorek, H.~Weinfurter, L.~Pezz\'e and A.~Smerzi,
Fisher information and multiparticle entanglement,
\emph{Phys. Rev. A} \textbf{85} (2012), 022321;
arXiv:1006.4366v3; doi:10.1103/PhysRevA.85.022321.
Section~III.A, Observation~1, equation~(9).
\bibitem{Toth78} G.~T\'oth,
Multipartite entanglement and high-precision metrology,
\emph{Phys. Rev. A} \textbf{85} (2012), 022322;
arXiv:1006.4368v3; doi:10.1103/PhysRevA.85.022322.
Observation~3, equation~(6).
\bibitem{CZY78} K.~Chen, Z.~Zhang and N.~Yu,
Optimal lower bound for quantum channel tomography in away-from-boundary regime,
arXiv:2601.10683v1 (15 January 2026), Theorems~1.1, 3.2 and~4.2.
\bibitem{CLL78} S.~Chen, J.~Li and A.~Liu,
An optimal tradeoff between entanglement and copy complexity for state tomography,
arXiv:2402.16353v1 (26 February 2024).
\bibitem{NZ78} A.~Nayak and X.~Zhou,
Optimal Low-Rank Quantum State Tomography with Bounded-Sample Joint Measurements,
arXiv:2609.10514v1 (9 September 2026), Theorems~1.1--1.2.

\bibitem{BPR80} S.~Basu, R.~Pollack and M.-F.~Roy,
\emph{Algorithms in Real Algebraic Geometry}, second edition,
Algorithms and Computation in Mathematics, vol.~10,
Springer, Berlin, Heidelberg, 2006;
doi:10.1007/3-540-33099-2.
Theorems~2.77 and~2.80 and Chapter~14.
\bibitem{PSTW80} A.~Pelecanos, J.~Spilecki, E.~Tang and J.~Wright,
Mixed state tomography reduces to pure state tomography,
arXiv:2511.15806v1 (19 November 2025).
Theorems~1.1--1.3 and~4.2, Section~3.3, Figures~8--9.

\bibitem{ZJ84} S.~Zhou and L.~Jiang,
Asymptotic theory of quantum channel estimation,
\emph{PRX Quantum} \textbf{2} (2021), 010343;
arXiv:2003.10559v3; doi:10.1103/PRXQuantum.2.010343.
\bibitem{KL84} E.~Knill and R.~Laflamme,
A theory of quantum error-correcting codes,
\emph{Phys. Rev. A} \textbf{55} (1997), 900--911;
arXiv:quant-ph/9604034; doi:10.1103/PhysRevA.55.900.
\bibitem{SKBG84} R.~Saini, J.~Kiukas, D.~Burgarth and A.~Gilchrist,
Characterizing Fisher information of quantum measurement,
arXiv:2512.15428v1 (2025).
\bibitem{MY84} P.~N.~Mayer and H.~Yun,
Kernel embedding for operator-valued measures and its application to
quantum tomography, arXiv:2605.25146v1 (2026).
\bibitem{YOPG84} S.~Yoshida, K.~Okigami, P.~M.~Posta and D.~Grinko,
Optimal learning of covariant quantum states and channels,
arXiv:2609.39280v1 (2026), Theorem~7, Corollary~9 and
Theorems~12--13, 15 and~18; full text consulted 5 October 2026.

\end{thebibliography}

\end{document}
